% Chapter 4. All quoted numbers come from chapters/ch04-numbers.tex, written by code/ch04/run_all.py.
% Generated by code/ch04/run_all.py. Do not edit by hand.
\newcommand{\amCallDiv}{8.1163}
\newcommand{\amCallDivPrem}{0.0156}
\newcommand{\amCallNoDiv}{10.4486}
\newcommand{\amPremium}{0.5165}
\newcommand{\amPremiumTree}{0.5175}
\newcommand{\amPutThousand}{6.0896}
\newcommand{\amPutTwoThousand}{6.0900}
\newcommand{\bandFeeGc}{0.17}
\newcommand{\bandFeeSp}{4.3}
\newcommand{\bandHi}{104.5505}
\newcommand{\bandLo}{103.8212}
\newcommand{\bandLoGc}{103.6448}
\newcommand{\bandLoSp}{99.4515}
\newcommand{\bandMid}{104.1852}
\newcommand{\bandProfit}{0.2495}
\newcommand{\bandQuote}{104.80}
\newcommand{\bandWidth}{0.73}
\newcommand{\bandWidthSp}{5.10}
\newcommand{\bndHalf}{83.62}
\newcommand{\bndHalfOne}{83.62}
\newcommand{\bndStart}{80.68}
\newcommand{\bndTenEnvT}{6.425}
\newcommand{\bndTenFirst}{73.26}
\newcommand{\bndTenFirstEnv}{72.23}
\newcommand{\bndTenFirstT}{9.890}
\newcommand{\bndTenSecond}{72.23}
\newcommand{\bndTenSecondT}{9.885}
\newcommand{\bndTenYear}{72.23}
\newcommand{\boundHi}{10.7843}
\newcommand{\boundLo}{1.9608}
\newcommand{\bsCall}{6.0401}
\newcommand{\bsDelta}{0.4496}
\newcommand{\bsPut}{5.5735}
\newcommand{\callQuote}{7.00}
\newcommand{\caseMonths}{309}
\newcommand{\cipBorrow}{1,370,000}
\newcommand{\cipFwd}{1.3535}
\newcommand{\cipOweCAD}{1,407,675}
\newcommand{\cipPoints}{\ensuremath{-}165}
\newcommand{\cipProfit}{6,725}
\newcommand{\cipQuote}{1.3600}
\newcommand{\cipReceive}{1,414,400}
\newcommand{\cipUSDEnd}{1,040,000}
\newcommand{\convNoY}{82.0252}
\newcommand{\convYield}{10.06}
\newcommand{\crrFifty}{6.0607}
\newcommand{\crrFiveHundred}{6.0422}
\newcommand{\crrHundred}{6.0532}
\newcommand{\crrHundredOne}{6.0311}
\newcommand{\crrNErrMax}{1.545}
\newcommand{\crrNErrMin}{\ensuremath{-}2.395}
\newcommand{\crrSlope}{\ensuremath{-}0.86}
\newcommand{\euCallDiv}{8.1007}
\newcommand{\euPutTwoThousand}{5.5725}
\newcommand{\exEightAm}{3.3218}
\newcommand{\exEightContD}{5.9804}
\newcommand{\exEightContU}{0.9804}
\newcommand{\exEightEu}{2.9220}
\newcommand{\exEightExD}{7.00}
\newcommand{\exFiveCall}{5.00}
\newcommand{\exFivePsiD}{0.50}
\newcommand{\exFivePsiU}{0.50}
\newcommand{\exFivePut}{7.50}
\newcommand{\exFourF}{1.3368}
\newcommand{\exFourGet}{1,373,500}
\newcommand{\exFourOwe}{1,370,250}
\newcommand{\exFourProfit}{3,250}
\newcommand{\exNineBasisIn}{2.8171}
\newcommand{\exNineBasisOut}{1.3606}
\newcommand{\exNineEx}{2.7399}
\newcommand{\exNineRoll}{\ensuremath{-}1.4565}
\newcommand{\exNineSpot}{4.1964}
\newcommand{\exOneDOne}{0.970874}
\newcommand{\exOneDTwo}{0.933352}
\newcommand{\exOneFwd}{4.020}
\newcommand{\exSevenDelta}{\ensuremath{-}0.50}
\newcommand{\exSevenP}{1.9608}
\newcommand{\exSevenQ}{0.60}
\newcommand{\exSixY}{11.88}
\newcommand{\exThreeV}{3.4940}
\newcommand{\exTwoF}{51.1378}
\newcommand{\fbEnd}{September 2026}
\newcommand{\fbFwdEnd}{September 2027}
\newcommand{\fbFwdStart}{January 2002}
\newcommand{\fbOneConst}{\ensuremath{-}0.19}
\newcommand{\fbOneHi}{1.02}
\newcommand{\fbOneLastDate}{September 2026}
\newcommand{\fbOneLo}{0.22}
\newcommand{\fbOneN}{297}
\newcommand{\fbOneRSq}{0.11}
\newcommand{\fbOneSE}{0.20}
\newcommand{\fbOneSlope}{0.62}
\newcommand{\fbOneTOne}{\ensuremath{-}1.88}
\newcommand{\fbOneTZero}{3.07}
\newcommand{\fbStart}{January 2001}
\newcommand{\fbTwoConst}{\ensuremath{-}0.37}
\newcommand{\fbTwoHi}{1.16}
\newcommand{\fbTwoLastDate}{September 2026}
\newcommand{\fbTwoLo}{0.11}
\newcommand{\fbTwoN}{285}
\newcommand{\fbTwoRSq}{0.09}
\newcommand{\fbTwoSE}{0.27}
\newcommand{\fbTwoSlope}{0.63}
\newcommand{\fbTwoTOne}{\ensuremath{-}1.38}
\newcommand{\fbTwoTZero}{2.36}
\newcommand{\ffFiveNeg}{\ensuremath{-}64.7}
\newcommand{\ffFivePos}{96.0}
\newcommand{\ffFiveZero}{15.6}
\newcommand{\ffMCFiveNeg}{\ensuremath{-}65.1}
\newcommand{\ffMCFiveNegSE}{0.6}
\newcommand{\ffMCFivePos}{95.8}
\newcommand{\ffMCFivePosSE}{0.6}
\newcommand{\ffMCFivePosZ}{0.2}
\newcommand{\ffOneNeg}{\ensuremath{-}4.3}
\newcommand{\ffOnePos}{4.8}
\newcommand{\ffOneZero}{0.3}
\newcommand{\ffPaths}{120,000}
\newcommand{\ffRhoStar}{\ensuremath{-}0.097}
\newcommand{\ffTenNeg}{\ensuremath{-}168.6}
\newcommand{\ffTenPos}{287.0}
\newcommand{\ffTenZero}{59.2}
\newcommand{\fraBuy}{959,693}
\newcommand{\fraRepay}{1,049,204}
\newcommand{\fraSell}{959,693}
\newcommand{\fwdAbove}{66}
\newcommand{\fwdErrMAE}{86}
\newcommand{\fwdErrMean}{27}
\newcommand{\fwdErrSE}{20}
\newcommand{\fwdErrT}{1.3}
\newcommand{\fwdIndex}{100.7528}
\newcommand{\fwdIndexLater}{103.3870}
\newcommand{\fwdValueAlt}{2.6079}
\newcommand{\fwdValueLater}{2.6079}
\newcommand{\gocDate}{29 September 2026}
\newcommand{\gocFwdOneTwo}{3.85}
\newcommand{\gocFwdOneTwoCont}{3.77}
\newcommand{\gocFwdSevenTen}{4.90}
\newcommand{\gocFwdSevenTenCont}{4.57}
\newcommand{\gocOneY}{2.97}
\newcommand{\gocSixM}{2.63}
\newcommand{\gocTenY}{3.99}
\newcommand{\gocThreeM}{2.39}
\newcommand{\gocTwoY}{3.37}
\newcommand{\gocZeroFive}{3.67}
\newcommand{\gocZeroTen}{4.00}
\newcommand{\hxBondFour}{97.9260}
\newcommand{\hxCheckTwo}{100.000000}
\newcommand{\hxDHalf}{0.980392}
\newcommand{\hxDOne}{0.959693}
\newcommand{\hxDOneHalf}{0.936920}
\newcommand{\hxDThree}{0.869504}
\newcommand{\hxDTwo}{0.914687}
\newcommand{\hxFwdInstTwo}{5.066}
\newcommand{\hxFwdOneTwo}{4.920}
\newcommand{\hxFwdTwoThree}{5.196}
\newcommand{\hxZHalf}{3.961}
\newcommand{\hxZOne}{4.114}
\newcommand{\hxZThree}{4.661}
\newcommand{\hxZTwo}{4.459}
\newcommand{\lrErrHundredOne}{\ensuremath{-}36.6}
\newcommand{\lrFiveHundredOne}{6.040087}
\newcommand{\lrHundredOne}{6.040051}
\newcommand{\lrSlope}{\ensuremath{-}1.99}
\newcommand{\naiveMAE}{79}
\newcommand{\opBond}{\ensuremath{-}57.1429}
\newcommand{\opDelta}{0.6667}
\newcommand{\opPrice}{9.5238}
\newcommand{\opPriceQ}{9.5238}
\newcommand{\opQ}{0.50}
\newcommand{\perpBoundary}{71.43}
\newcommand{\psiDown}{0.2520}
\newcommand{\psiMid}{0.3784}
\newcommand{\psiSum}{0.9804}
\newcommand{\psiUp}{0.3500}
\newcommand{\putFair}{11.3431}
\newcommand{\putGap}{1.6569}
\newcommand{\putQuote}{13.00}
\newcommand{\qUpMax}{55}
\newcommand{\qUpMin}{10}
\newcommand{\qrnDown}{25.70}
\newcommand{\qrnMid}{38.60}
\newcommand{\qrnUp}{35.70}
\newcommand{\repBond}{107.8431}
\newcommand{\repCall}{0.5000}
\newcommand{\repStock}{\ensuremath{-}1.0000}
\newcommand{\rollApproxMean}{3.68}
\newcommand{\rollBasisCorr}{\ensuremath{-}0.90}
\newcommand{\rollContango}{54.0}
\newcommand{\rollConvTerm}{\ensuremath{-}4.55}
\newcommand{\rollExAnn}{\ensuremath{-}0.02}
\newcommand{\rollGapMaxBp}{5.9}
\newcommand{\rollGapMedBp}{1.1}
\newcommand{\rollMCContango}{49.2}
\newcommand{\rollMCHi}{8.4}
\newcommand{\rollMCLo}{\ensuremath{-}7.1}
\newcommand{\rollMCMean}{0.68}
\newcommand{\rollMCSE}{0.05}
\newcommand{\rollMCSd}{4.73}
\newcommand{\rollMCSpotHi}{1.3}
\newcommand{\rollMCSpotLo}{\ensuremath{-}1.3}
\newcommand{\rollPaths}{10,000}
\newcommand{\rollPremTerm}{5.17}
\newcommand{\rollRate}{3}
\newcommand{\rollRollAnn}{0.65}
\newcommand{\rollRollSE}{1.20}
\newcommand{\rollSpotAnn}{\ensuremath{-}0.67}
\newcommand{\rollSpotAnnAbs}{0.67}
\newcommand{\rollSpotFortyFall}{24}
\newcommand{\rollTheory}{0.62}
\newcommand{\rollTotalAnn}{2.98}
\newcommand{\rollWealthAnn}{2.99}
\newcommand{\rollWealthMean}{3.69}
\newcommand{\rollYears}{40}
\newcommand{\schFinf}{71.81}
\newcommand{\schOneHigh}{79.77}
\newcommand{\schOneLow}{64.74}
\newcommand{\schSFlat}{73.67}
\newcommand{\schSStar}{71.81}
\newcommand{\snipAm}{6.0896}
\newcommand{\snipEu}{5.5715}
\newcommand{\sprMean}{23}
\newcommand{\sprPos}{70}
\newcommand{\sprSD}{61}
\newcommand{\tpAmerican}{7.1429}
\newcommand{\tpContDown}{11.4286}
\newcommand{\tpEuropean}{5.4422}
\newcommand{\tpExDown}{15.00}
\newcommand{\tpPremium}{1.7007}
\newcommand{\treeDelta}{0.4496}

\begin{refsection}
\setcounter{chapter}{3}
\chapter{Arbitrage, Forwards, Futures and Binomial Trees}
\label{ch:arbitrage}

A farmland fund that agrees today to sell next year's crop, a pension plan that rolls commodity futures
for inflation protection, and a manager who hedges the currency of a foreign property all depend on one
question: what should a price for future delivery be? This chapter answers it without forecasting
anything. Two portfolios with the same future cash flows must cost the same today; otherwise someone
buys the cheap one, sells the dear one and keeps the difference. That single rule fixes forward
prices, links interest rates across currencies and, applied one step at a time on a tree, prices
options.

We start with what ruling out arbitrage assumes and what it costs to enforce (\cref{sec4:why}),
then use the rule to turn bill and bond yields into a discount curve (\cref{sec4:curve}). In a
general one-period market the same rule becomes the fundamental theorem of asset pricing
(\cref{sec4:ftap}), and it goes on to fix the prices of forwards and futures on assets that pay
income or cost money to store (\cref{sec4:fwd,sec4:fut,sec4:roll}). On a binomial tree it prices options by backward induction, converges to the Black-Scholes formula
as the steps shrink and handles early exercise
(\cref{sec4:binom,sec4:limit,sec4:american}). The case study then asks whether a
Government of Canada forward rate, which arbitrage fixes as a price, also tells us where rates will go
(\cref{sec4:case}).

\begin{keybox}[How to read this chapter]
The numbers in the text are computed by seven scripts in \texttt{code/ch04/}
(\texttt{term\_structure}, \texttt{no\_arbitrage}, \texttt{forwards\_carry}, \texttt{commodity\_roll},
\texttt{binomial\_trees}, \texttt{case\_forward\_rates} and \texttt{exercises}). The tests in
\texttt{code/ch04/test\_*.py} check them against closed forms, enumeration and simulation, an
independent R script recomputes the curve, the trees and the case study, and GNU Octave checks the
price bounds with its own linear programming solver and the trees by backward induction. The chapter proves its main finite-state results and names the standard results from linear
programming and probability that it uses without proof. Market data and published results are cited where they appear. The simulations
draw from fixed random streams of the seed 404, so they reproduce exactly in the recorded
software environment.

Notation is defined where first used; the reference box at the end of \cref{sec4:why} records
symbols reused with different meanings in later sections.
\end{keybox}

% =====================================================================================
\section{Prices that cannot be undercut}
\label{sec4:why}

An \emph{arbitrage} is a trade that requires no positive initial investment, can never lose, and
either brings in money today or sometimes gains later. The definition is strict on purpose. A trade that is expected to make money but can
lose is a bet, however good, and a trade that cannot lose but costs something today is an investment.
Only when both conditions hold does the trade make money from nothing, and traders who take it move
prices until the opportunity is gone. Pricing by no arbitrage assumes that this has already happened.

The assumption is weaker than market efficiency in the sense of \textcite{fama1970}. It says nothing
about whether prices reflect information about the future. It only says that prices are consistent
with each other: a forward cannot be quoted far from the spot price carried forward, and a call
cannot be worth more than the stock it is written on. These conditions hold whatever anyone believes
about the future, up to the cost of trading on them.

Enforcing them is not free. Each arbitrage argument below needs someone who can borrow, lend, sell
short and store, and in practice each of those has a cost. The result is a band of prices rather than
a single price, and \cref{sec4:fwd} measures how wide the band is. For an alternative investment the
band can be wide: nobody sells a timber tract short to arbitrage a mispriced lumber forward.

\begin{keybox}[Symbols reused between sections]
Within a section each letter keeps one meaning, apart from running indices ($j$, $k$, $m$) and
integration variables such as the $u$ in $\Lambda(u,T)$. The letters whose meaning changes between
sections are noted here, and they are the ones a reader is likely to confuse. $X$ is a claim's payoff in \cref{sec4:ftap} and the log spot price in
\cref{sec4:roll}; $V$ is the value of a forward in \cref{sec4:fwd}, the investor's wealth in
\cref{sec4:roll} and an option value on a tree; $m$ is the compounding frequency in \cref{sec4:curve},
the deterministic part $m(t)$ of the short rate in \cref{sec4:fut} and the number of rolls in
\cref{sec4:roll}; $c$ is a coupon rate in \cref{sec4:curve} and $c_k$ a European call value in
\cref{sec4:american}; $w$ is an interpolation weight in \cref{sec4:curve} and $w_{k,j}=V_{k,j}+S_{k,j}$ in
\cref{sec4:american}; and in \cref{sec4:limit} $d_1$ and $d_2$ are the Black-Scholes arguments, not down
factors. $S$ is a price, of an asset or at a tree node; the number of states in \cref{sec4:ftap} is $J$.
$P$ is a bond or put price, and the real-world measure is $\mathbb P$. In \cref{sec4:case}, $a$ and
$b$ are regression coefficients; in \cref{sec4:american}, $b_k$ and $b$ are exercise levels; and a
superscript $a$ or $b$ marks an ask or a bid. The rate $r$ is constant except in the Vasicek model
of \cref{sec4:fut}, where $r_t$ is the short rate. On a
tree $u$ and $d$ are the up and down factors, $R$ the gross interest per step and $h$ the step
length; in \cref{sec4:ftap}, $h$ is instead the upper price bound, and in \cref{sec4:roll} the roll
step. In \cref{sec4:fut}, $\Lambda(u,T)=(1-e^{-\kappa(T-u)})/\kappa$ is a zero coupon bond's sensitivity
to the short rate, so its volatility is $\sigma_r\Lambda$; it is unrelated to the storage rate
$\lambda$. In \cref{sec4:roll}, $b_t(T)$ is the basis and $g_k$ a log futures change. $g$ is a payoff function in \cref{sec4:binom} and an exercise payoff in \cref{sec4:american}; $n$ is an
asset index in \cref{sec4:ftap}, the number of settlement dates in \cref{sec4:fut} and otherwise the
number of tree steps; $\epsilon$ is a small perturbation in \cref{sec4:ftap}, $\varepsilon_h$ a log step in
\cref{sec4:limit} and $\varepsilon_{t+12}$ a regression error in \cref{sec4:case}. $\tau$ is a term or a time to expiry, except in
\cref{sec4:american}, where it is a stopping time. Continuously compounded rates are $r$ and simple money market rates are
$i$ or $y$. An asset's income yield is $\delta$, a commodity's storage rate $\lambda$ and its implied
convenience yield $\eta$. $K$ is a strike or delivery price, $k$ the contract rate of \cref{prop4:fra} and otherwise an index, $\theta$ a
portfolio, $q$ a risk-neutral probability and $B$ the money market account. $F_0$ is a forward price, $F(t_1,t_2)$ a forward
interest rate, and $G$ a futures price.
\end{keybox}

% =====================================================================================
\section{Discount factors, zero rates and forward rates}
\label{sec4:curve}

\begin{definition}\label{def4:curve}
The \emph{discount factor} $D(t)$ is the price today of one dollar paid with certainty at time $t$,
with $D(0)=1$. The continuously compounded \emph{zero rate} is $z(t)=-\ln D(t)/t$. For $t_1<t_2$, the
simple \emph{forward rate} for the period $[t_1,t_2]$ and the \emph{instantaneous forward rate} at $t$
are
\begin{equation}\label{eq4:fwd}
  F(t_1,t_2)=\frac{1}{t_2-t_1}\Bigl(\frac{D(t_1)}{D(t_2)}-1\Bigr),\qquad
  f(t)=-\frac{d}{dt}\ln D(t).
\end{equation}
\end{definition}

A money market rate $y$ quoted on a simple basis for a term $t$ means $D(t)=1/(1+yt)$; a rate $y$
compounded $m$ times a year means $D(t)=(1+y/m)^{-mt}$. The forward rate is a price, not a forecast,
as the next result shows: anyone can lock it in today with two bonds.

\begin{proposition}[Locking in a forward rate]\label{prop4:fra}
Suppose zero-coupon bonds maturing at $t_1$ and $t_2$ trade at $D(t_1)$ and $D(t_2)$ and there is no
arbitrage. Then an agreement made today to borrow $1$ at $t_1$ and repay $1+k(t_2-t_1)$ at $t_2$ has
zero value if and only if $k=F(t_1,t_2)$. Also,
\begin{equation}\label{eq4:zeroavg}
  D(t)=\exp\Bigl(-\int_0^t f(s)\,ds\Bigr),\qquad z(t)=\frac1t\int_0^t f(s)\,ds ,
\end{equation}
so the zero rate is the average of the instantaneous forward rates up to $t$.
\end{proposition}

\begin{proof}
The agreement pays $+1$ at $t_1$ and $-(1+k\tau)$ at $t_2$, with $\tau=t_2-t_1$. Buying one bond
maturing at $t_1$ and selling $1+k\tau$ bonds maturing at $t_2$ produces exactly these cash flows.
Two portfolios with the same payoffs must have the same price when there is no arbitrage; this is the
law of one price, proved as \cref{prop4:lop} in the next section. So the agreement is worth
$D(t_1)-(1+k\tau)D(t_2)$. This is zero exactly when $1+k\tau=D(t_1)/D(t_2)$, which is
\cref{eq4:fwd}. For \cref{eq4:zeroavg}, integrate $f=-(\ln D)'$ from $0$ to $t$ and use $\ln D(0)=0$.
\end{proof}

\paragraph{Bootstrapping a curve} Markets quote bill yields and coupon bond yields, and the discount
factors have to be inferred from them. Each quote is converted using the instrument's actual cash
flows, quotation basis and settlement convention. In the simplest case, a coupon bond paying $c/m$ at $t_k=k/m$, $k=1,\dots,mT$, and $1$ at $T$ is worth
$\sum_k (c/m)D(t_k)+D(T)$. A \emph{par yield} is the coupon that makes this equal to $1$. Given
discount factors at a set of knots, we fill the gaps by interpolating $\ln D$ linearly, which is the
same as holding the instantaneous forward rate constant between knots.

\begin{proposition}[Bootstrapping]\label{prop4:boot}
Let discount factors be known at knots $0=s_0<s_1<\dots<s_j$ and let a bond with coupon $c>0$ mature
at $T>s_j$. Suppose the coupons paid up to $s_j$ are worth less than par, $(c/m)\sum_{t_k\le s_j}D(t_k)<1$,
an assumption the existence argument needs and every bond below satisfies. With log-linear interpolation between $s_j$ and $T$, there is
exactly one $D(T)\in(0,\infty)$ that prices the bond at par.
\end{proposition}

\begin{proof}
For $t\in(s_j,T]$ the interpolated factor is $D(t)=D(s_j)^{1-w}D(T)^{w}$ with $w=(t-s_j)/(T-s_j)\in(0,1]$,
which is continuous and strictly increasing in $D(T)$. Coupon dates up to $s_j$ do not depend on
$D(T)$. Hence the bond price $P(D_T)$ is continuous and strictly increasing. It tends to
$(c/m)\sum_{t_k\le s_j}D(t_k)<1$ as $D_T\to0$ and to $\infty$ as $D_T\to\infty$. The intermediate value
theorem gives a root, and strict monotonicity makes it unique.
\end{proof}

\begin{example}{Building a curve by hand}\label{ex4:boot}
The instruments here are idealised: bill yields are simple rates for exact terms, and each bond pays
semiannual coupons at 0.5, 1, 1.5, \ldots{} years and prices at par. Bills for six months and one year yield $4.00\%$ and $4.20\%$ on a simple basis, so $D(0.5)=\hxDHalf$
and $D(1)=\hxDOne$. A two year bond with semiannual coupons has a par yield of $4.50\%$. The unknown
$D(2)$ enters the price directly and through the interpolated value
\[
  D(1.5)=\sqrt{D(1)D(2)};
\]
solving gives
$D(2)=\hxDTwo$ and $D(1.5)=\hxDOneHalf$. A three year par yield of $4.70\%$ then gives
$D(3)=\hxDThree$. The zero rates rise from $\hxZHalf\%$ at six months to $\hxZThree\%$ at three
years, and in this example the forward rates lie above them: $F(1,2)=\hxFwdOneTwo\%$ and
$F(2,3)=\hxFwdTwoThree\%$, both simple rates. Compared on the same compounding basis, a forward rate
exceeds the zero rate at the end of its period whenever the zero rate rises over that period; a rising
zero curve need not put every forward above every earlier zero rate. A three year bond with an annual coupon of
$4\%$ is worth \hxBondFour{} per 100. To lock in $F(1,2)$ on 1,000,000, buy 1,000,000 face of the
one year bond for \fraBuy{} and sell \fraRepay{} face of a two year zero coupon bond (a strip), which raises \fraSell:
the position costs nothing, pays 1,000,000 in a year and owes \fraRepay{} in two.
\end{example}

\Cref{fig4:goc} applies the same method to the Government of Canada curve on \gocDate, the last
business day of September (30 September, the National Day for Truth and Reconciliation, is a Bank
holiday with no published rates). Bills for three
months, six months and one year yield \gocThreeM\%, \gocSixM\% and \gocOneY\%; benchmark bonds for two
to ten years yield from \gocTwoY\% to \gocTenY\%. Two approximations are made. Bill yields are treated
as simple rates on an actual/365 basis with exact terms of a quarter, half and whole year, and each
benchmark bond is treated as a par bond with semiannual coupons maturing exactly at its nominal term.
Benchmark bonds trade close to par but rarely exactly at it, so this is a transparent bootstrap of
the quoted benchmark yields, showing the shape they imply, and not a curve fitted to the prices of
actual bonds. The bootstrap reprices these idealised inputs to within $10^{-15}$, which checks
the bootstrap; it says nothing about how well the curve prices actual bonds with their own coupons and
dates. The five and ten year zero rates are \gocZeroFive\% and \gocZeroTen\%. The simple forward rate
$F(1,2)$ is \gocFwdOneTwo\% and $F(7,10)$ is \gocFwdSevenTen\% a year (\gocFwdOneTwoCont\% and
\gocFwdSevenTenCont\% continuously compounded, the levels in \cref{fig4:goc}).

\begin{figure}[t]
  \centering
  \includegraphics[width=0.92\textwidth]{ch04/fig01_goc_curve}
  \caption{Zero and instantaneous forward rates bootstrapped from Government of Canada bill and
  benchmark bond yields on \gocDate, with discount factors interpolated log-linearly; both curves are
  continuously compounded. The forward rate is the slope of $-\ln D$, so it jumps at each knot and lies
  above the zero rate wherever the zero rate rises. The steps come from the interpolation; they are
  not separately observed market quotes.}
  \label{fig4:goc}
\end{figure}

Monotone-convex interpolation or a smoothing spline removes the steps, at the cost of fitting extra
parameters. \Cref{eq4:zeroavg} holds for every one of these curves, the log-linear one
included, because it needs only that $\ln D$ be the integral of $f$; what smoothing changes is the
forward rates between the observed maturities.

% =====================================================================================
\section{Arbitrage and state prices}
\label{sec4:ftap}

\begin{definition}\label{def4:market}
A \emph{one-period market} has $J$ states of the world at time $1$ and $N$ assets. Asset $n$ costs
$p_n$ today and pays $A_{sn}$ in state $s$, so the payoff matrix is $A\in\mathbb R^{J\times N}$. A
portfolio $\theta\in\mathbb R^N$ costs $p^\top\theta$ and pays $A\theta$. It is an \emph{arbitrage}
if $p^\top\theta\le0$ and $A\theta\ge0$ in every state, with $p^\top\theta<0$ or $A\theta\ne0$.
A \emph{state price vector} is a $\psi\in\mathbb R^J$ with $\psi_s>0$ for every $s$ and
$A^\top\psi=p$.
\end{definition}

The state price $\psi_s$ is what the market charges today for one dollar paid in state $s$ and nothing
otherwise. If a riskless bond paying $R$ in every state costs $1$, then $\sum_s\psi_s=1/R$, and
$q_s=R\psi_s$ are positive numbers that add to one. They are called \emph{risk-neutral probabilities}
because every price becomes a discounted expectation under them:
\begin{equation}\label{eq4:rn}
  p_n=\sum_s\psi_sA_{sn}=\frac1R\sum_sq_sA_{sn}=\frac1R\,\E_Q[A_n].
\end{equation}

\begin{proposition}[Law of one price]\label{prop4:lop}
If there is no arbitrage and $A\theta=A\theta'$, then $p^\top\theta=p^\top\theta'$.
\end{proposition}

\begin{proof}
If $p^\top\theta<p^\top\theta'$, the portfolio $\theta-\theta'$ pays $0$ in every state and has
negative cost, which is an arbitrage. The other inequality is ruled out the same way.
\end{proof}

The converse half of the next theorem rests on one geometric fact: if a subspace and a compact convex
set do not meet, some hyperplane through the origin contains the subspace and leaves the whole set
strictly on one side. Orthogonal projection gives the proof.

\begin{lemma}[Separation]\label{lem4:sep}
Let $M\subset\mathbb R^k$ be a linear subspace and $\mathcal C\subset\mathbb R^k$ a nonempty compact convex
set with $M\cap\mathcal C=\emptyset$. Then there is a $v\in\mathbb R^k$ with $v^\top m=0$ for every $m\in M$ and
$v^\top x>0$ for every $x\in\mathcal C$.
\end{lemma}

\begin{proof}
For $x\in\mathcal C$ let $\pi(x)$ be the orthogonal projection of $x$ onto $M$, so that $\lVert x-\pi(x)\rVert$
is the distance from $x$ to $M$. This function is continuous, so it attains its minimum on the compact
set $\mathcal C$ at some
$x^*$; write $m^*=\pi(x^*)$ and $v=x^*-m^*$. Then $v\ne0$ because $x^*\notin M$, and $v$ is orthogonal
to $M$ by the definition of the projection, which is the first claim. For the second, take any
$x\in\mathcal C$ and $t\in(0,1]$. The point $x^*+t(x-x^*)$ lies in $\mathcal C$ by convexity, so by the
choice of $x^*$ its distance to $M$ is at least $\lVert v\rVert$. That distance is at most its distance
to the point $m^*$ of $M$, which is $\lVert v+t(x-x^*)\rVert$, so
\begin{equation*}
  \lVert v\rVert^2\le\lVert v+t(x-x^*)\rVert^2=\lVert v\rVert^2+2t\,v^\top(x-x^*)+t^2\lVert x-x^*\rVert^2 .
\end{equation*}
Dividing by $t$ and letting $t\to0$ gives $v^\top x\ge v^\top x^*=v^\top(v+m^*)=\lVert v\rVert^2>0$.
\end{proof}

\begin{theorem}[Fundamental theorem of asset pricing, one period]\label{thm4:ftap}
A one-period market has no arbitrage if and only if a state price vector exists.
\end{theorem}

\begin{proof}
Suppose $\psi$ is a state price vector and $A\theta\ge0$. Then
$p^\top\theta=\psi^\top A\theta\ge0$, with equality only if $A\theta=0$ because every $\psi_s>0$. So
a portfolio with $A\theta\ge0$ and $p^\top\theta\le0$ has $p^\top\theta=0$ and $A\theta=0$, and is not
an arbitrage.

Conversely, suppose there is no arbitrage. Consider the subspace of $\mathbb R^{1+J}$ of cost and payoff
pairs, $M=\{(-p^\top\theta,\,A\theta):\theta\in\mathbb R^N\}$, and the simplex
$\mathcal C=\{x\in\mathbb R^{1+J}:x\ge0,\ \sum_ix_i=1\}$. A point of $M\cap\mathcal C$ would be a portfolio with
$-p^\top\theta\ge0$, $A\theta\ge0$ and not all of these zero, which is an arbitrage. So $M\cap\mathcal C$ is
empty, and \cref{lem4:sep} gives $v=(v_0,v_1,\dots,v_J)$ with $v^\top x>0$ on $\mathcal C$ and $v^\top m=0$ on
$M$. Applying the first property to the corners of the simplex shows $v_i>0$ for every $i$. The
second says $-v_0p^\top\theta+\sum_sv_s(A\theta)_s=0$ for every $\theta$, that is
$A^\top(v_{1:J}/v_0)=p$. So $\psi_s=v_s/v_0$ is a state price vector.
\end{proof}

The one-period result goes back to \textcite{ross1978}. \textcite{harrison1979} tied the absence
of arbitrage to a martingale measure in a general multiperiod setting, through the viability of prices
as an economic equilibrium, and \textcite{harrison1981} proved the exact equivalence for markets with
finitely many states and dates and developed the theory of continuous trading. The finite-state
formulation and proof given here follow \textcite{follmer2016} and \textcite{duffie2001}. Chapter~5 uses the
corollary that follows to read option prices as expectations, and Chapter~6 uses the bounds of the
proposition after it for claims that cannot be replicated.

\begin{corollary}\label{cor4:complete}
Suppose there is no arbitrage. The state price vector is unique if and only if the market is
\emph{complete}, meaning that $A$ has rank $J$, so that every payoff in $\mathbb R^J$ is the payoff of some
portfolio. In a complete market every claim $X$ has the unique price $\psi^\top X$, which is the cost
of the portfolio that replicates it.
\end{corollary}

\begin{proof}
State price vectors are the positive solutions of $A^\top\psi=p$. If $A$ has rank $J$, then $A^\top$
is one to one and there is at most one solution. If the rank is less than $J$, there is a nonzero
$\zeta$ with $A^\top\zeta=0$, and $\psi+\epsilon\zeta$ is a second state price vector for small
$\epsilon$. When $A\theta=X$, $\psi^\top X=\psi^\top A\theta=p^\top\theta$.
\end{proof}

\begin{proposition}[Price bounds in an incomplete market]\label{prop4:bounds}
Suppose there is no arbitrage, that some portfolio has a strictly positive payoff in every
state, and let $X\in\mathbb R^J$ be any claim. The prices of $X$ that can be added to the market
without creating an arbitrage are the values $\psi^\top X$ over strictly positive state price vectors
$\psi$. They form an interval whose closure is $[\ell,h]$, where
\begin{equation}\label{eq4:bounds}
  h=\max\{\psi^\top X:\psi\ge0,\ A^\top\psi=p\}=\min\{p^\top\theta:A\theta\ge X\},
\end{equation}
and $\ell$ is the corresponding minimum, equal to $\max\{p^\top\theta:A\theta\le X\}$. The upper
bound $h$ is the cost of the cheapest portfolio that \emph{super-replicates} $X$, and $\ell$ is the
value of the dearest portfolio that sub-replicates it.
\end{proposition}

\begin{proof}
By \cref{thm4:ftap} applied to the enlarged market, a price $x$ for $X$ is free of arbitrage exactly
when some strictly positive $\psi$ satisfies $A^\top\psi=p$ and $\psi^\top X=x$. The set of such
$\psi$ is convex, so the set of such $x$ is an interval. Its closure is obtained by allowing
$\psi\ge0$: any $\psi\ge0$ with $A^\top\psi=p$ is the limit of $(1-\epsilon)\psi+\epsilon\psi^0$,
where $\psi^0$ is strictly positive. The set $\{\psi\ge0:A^\top\psi=p\}$ is bounded, because a portfolio with a strictly positive
payoff has a positive price that caps the total mass of any such $\psi$; it is also closed, hence compact, so the relaxed problem attains both endpoints. Whether an endpoint
is itself an arbitrage-free price is taken up after the example. The two expressions for $h$ are a primal and dual pair
of linear programs. Weak duality is
direct: if $\psi\ge0$, $A^\top\psi=p$ and $A\theta\ge X$, then
$p^\top\theta=\psi^\top A\theta\ge\psi^\top X$. Equality of the two optimal values is the strong
duality theorem of linear programming, which applies because both programs are feasible (the first by
\cref{thm4:ftap}, the second because a large enough holding of the portfolio with a strictly positive
payoff dominates $X$).
\end{proof}

\begin{example}{State prices in three states}\label{ex4:states}
A stock priced at 100 moves to 120, 100 or 80, a bond pays $R=1.02$ for a price of 1, and a call on the
stock struck at 100 is quoted at \callQuote. With three assets and three states the market is complete,
and solving $A^\top\psi=p$ gives $\psi=(\psiUp,\psiMid,\psiDown)$. All three are positive, so there is no
arbitrage, and they add to $\psiSum=1/1.02$. The risk-neutral probabilities are \qrnUp\%, \qrnMid\% and
\qrnDown\%. Now a dealer bids \putQuote{} for a put struck at 110, which pays $0$, $10$ and $30$. The
state prices value it at \putFair. The replicating portfolio holds \repBond{} bonds and \repCall{} calls and
is short one share, and costs \putFair; selling the put at \putQuote{} and buying the replica
locks in \putGap{} today with no exposure at all. A linear feasibility check based on
\cref{thm4:ftap} finds an arbitrage in the four-asset market and none in the three-asset one.
\end{example}

Remove the call from that market and it becomes incomplete: two assets cannot span three states. The
risk-neutral probabilities must still price the bond and the stock, which forces $q_u-q_d=0.1$ and
leaves $q_u$ free between \qUpMin\% and \qUpMax\%. \Cref{fig4:bounds} shows the resulting range of call
prices for every strike. At $K=100$ the closure of the range is $[\boundLo,\boundHi]$. The upper end is
the cost of half a share and a loan of $40/1.02$, which pays 20, 10 and 0, at least as much as the call
in every state. The lower end is the value of one share and a loan of $100/1.02$, which pays 20, 0 and
$-20$, never more than the call. Neither portfolio replicates the call, which is not attainable with
two assets, so neither endpoint is itself an arbitrage-free price: selling the call at \boundHi{} and
buying the super-replicating portfolio would lock in a gain in the middle state at no cost. The
arbitrage-free prices are the open interval $(\boundLo,\boundHi)$, as \textcite{follmer2016} show in
general for a claim that cannot be replicated. The quoted price of \callQuote{} sits inside, as it
must. Linear programming solves the primal and dual problems separately, and they agree to $10^{-13}$.

\begin{figure}[t]
  \centering
  \includegraphics[width=0.86\textwidth]{ch04/fig02_bounds}
  \caption{No-arbitrage prices of a call when only the stock and a bond trade and the stock can end at
  120, 100 or 80. The band is the closure of the values $\psi^\top X$ over strictly positive state price vectors.
  Its edges, the super-replication cost and the sub-replication value of \cref{prop4:bounds}, are
  not arbitrage-free prices themselves, except at $K\le80$ and $K\ge120$, where the call can be
  replicated and the band closes to a point.}
  \label{fig4:bounds}
\end{figure}

% =====================================================================================
\section{Forward prices by cash and carry}
\label{sec4:fwd}

\begin{definition}\label{def4:forward}
A \emph{forward contract} agreed at time $0$ obliges the long side to buy one unit of an asset at time
$T$ for the delivery price $K$; its payoff at $T$ is $S_T-K$. The \emph{forward price} $F_0$ is the
delivery price that gives the contract zero value at inception.
\end{definition}

In this section the interest rate $r$ is constant and continuously compounded. An asset may pay income
at a continuous yield $\delta$, as an equity index pays dividends or a foreign currency earns interest.

\begin{proposition}[Cost of carry]\label{prop4:carry}
If there is no arbitrage, the forward price of an asset with income yield $\delta$ is
\begin{equation}\label{eq4:carry}
  F_0=S_0e^{(r-\delta)T},
\end{equation}
and at a later time $t<T$ a long forward with delivery price $K$ is worth
\begin{equation}\label{eq4:fwdvalue}
  V_t=S_te^{-\delta(T-t)}-Ke^{-r(T-t)}=(F_t-K)e^{-r(T-t)} .
\end{equation}
\end{proposition}

\begin{proof}
Buy $e^{-\delta(T-t)}$ units of the asset at $t$ and reinvest the income in the asset; the holding grows to
exactly one unit at $T$. Borrow $Ke^{-r(T-t)}$ at the same time, so that $K$ is owed at $T$. The
portfolio pays $S_T-K$ at $T$, which is the forward payoff, and costs
$S_te^{-\delta(T-t)}-Ke^{-r(T-t)}$. By the law of one price this is the value of the forward. Setting it
to zero at $t=0$ gives \cref{eq4:carry}, and substituting $F_t=S_te^{(r-\delta)(T-t)}$ gives the second form
of \cref{eq4:fwdvalue}.
\end{proof}

The proof is a recipe for the arbitrage. If a forward is quoted above $S_0e^{(r-\delta)T}$, buy the
asset with borrowed money and sell the forward: this is \emph{cash and carry}. If it is quoted below,
sell the asset short, lend the proceeds and buy the forward: this is \emph{reverse cash and carry}. The
reverse trade needs someone willing to lend the asset, and the short seller must pass on to the lender
the income the asset pays while the short is open. For an index at 100 with $r=4\%$ and $\delta=2.5\%$,
the six month forward is \fwdIndex. Three months later, with the
index at 103, the new forward price is \fwdIndexLater{} and the old contract is worth \fwdValueLater,
whichever form of \cref{eq4:fwdvalue} is used.

\paragraph{Storage and convenience} A physical commodity costs money to store and pays no income. If the
commodity can be bought now and stored until $T$, with storage costs at a continuous rate $\lambda$,
cash and carry gives the upper bound $F_0\le S_0e^{(r+\lambda)T}$. The reverse trade needs someone to
sell the commodity now and buy it back later, and an owner who needs the stock to keep a refinery or a
mill running will not do that for a small gain. Holding inventory gives a benefit that a forward
position does not, and \textcite{kaldor1939} called it the \emph{convenience yield}. It is inferred
from prices; nobody receives it as a cash payment.

\begin{definition}\label{def4:convenience}
The \emph{implied convenience yield} $\eta$ of a forward on a commodity with storage rate $\lambda$ is
defined by $F_0=S_0e^{(r+\lambda-\eta)T}$, that is $\eta=r+\lambda-\ln(F_0/S_0)/T$.
\end{definition}

The upper bound says $\eta\ge0$; nothing bounds $\eta$ from above. With the spot at 80, $r=4\%$ and
$\lambda=1\%$, a six month forward at 78 implies $\eta=\convYield\%$ a year. Without a convenience yield
the forward would be \convNoY. A high implied convenience yield is consistent with scarce inventory,
though other frictions can affect the inference, and \cref{sec4:roll} turns this into a model.

\begin{proposition}[No-arbitrage band with frictions]\label{prop4:band}
Suppose trading is unrestricted in size, income is paid as a continuous yield, there are no margin,
collateral or other costs beyond those named, and any quantity of the asset can be bought at the ask $S^a$ and sold short at the bid
$S^b\le S^a$, money can be borrowed at $r_b$ and lent at $r_\ell\le r_b$, short sale proceeds earn
$r_\ell$, and a short seller pays the asset's income yield $\delta$ to the lender. Then there is no
arbitrage between the asset, the money market and a forward if and only if
\begin{equation}\label{eq4:band}
  S^be^{(r_\ell-\delta)T}\le F_0\le S^ae^{(r_b-\delta)T}.
\end{equation}
\end{proposition}

\begin{proof}
If $F_0>S^ae^{(r_b-\delta)T}$, borrow $S^ae^{-\delta T}$ at $r_b$, buy $e^{-\delta T}$ units at the ask,
reinvest the income in the asset and sell the forward. At $T$ deliver the unit, receive $F_0$ and repay
$S^ae^{(r_b-\delta)T}$, keeping the positive difference. If $F_0<S^be^{(r_\ell-\delta)T}$, sell
$e^{-\delta T}$ units short at the bid, lend the proceeds at $r_\ell$ and buy the forward; the income
owed on the short is met by shorting further units at the same yield, so the short grows to one unit
at $T$. The loan then returns $S^be^{(r_\ell-\delta)T}$, the forward delivers a unit for $F_0$, and the
short is closed with a profit.

Conversely, suppose \cref{eq4:band} holds. The function $(S,r)\mapsto Se^{(r-\delta)T}$ is continuous on
$[S^b,S^a]\times[r_\ell,r_b]$ and takes the two ends of the band at opposite corners, so some $S$ and $r$
in those ranges satisfy $F_0=Se^{(r-\delta)T}$. In a frictionless market with spot $S$ and rate $r$ this
forward price admits no arbitrage by \cref{prop4:carry}. Under the linear trading terms of the proposition, take any trade available with frictions, in which, without loss of
generality, the asset is either bought or sold short and money is either borrowed or lent: it
buys at prices of at least $S$, sells at no more than $S$, borrows at rates of at least $r$ and lends
at no more than $r$. So it costs at least as much today and pays no more at $T$ than the same trade at
the frictionless prices $S$ and $r$. If it were an arbitrage, the frictionless trade would be one too,
which is impossible.
\end{proof}

\Cref{fig4:band} draws the band for a spot quoted at 99.95 bid and 100.05 ask, lending at 3.8\% and
borrowing at 4.4\%. For one year it runs from \bandLo{} to \bandHi, around a frictionless price of
\bandMid{} at the midpoint rate. A quote of \bandQuote{} lies above the band, and cash and carry earns
\bandProfit{} per unit at expiry. The width grows with maturity because the gap between borrowing and
lending rates compounds.

\begin{figure}[t]
  \centering
  \includegraphics[width=0.86\textwidth]{ch04/fig03_band}
  \caption{The no-arbitrage band of \cref{prop4:band} for a spot at 99.95 bid and 100.05 ask, lending at
  3.8\% and borrowing at 4.4\%, with no income. A forward quoted above the upper line can be sold
  against a purchase of the asset with borrowed money.}
  \label{fig4:band}
\end{figure}

\paragraph{Covered interest parity} A currency is an asset whose income is the foreign interest rate.
With the exchange rate $S$ quoted as domestic currency per unit of foreign currency and simple money
market rates $i_d$ and $i_f$ for a term $\tau$, the argument of \cref{prop4:carry} gives
\begin{equation}\label{eq4:cip}
  F=S\,\frac{1+i_d\tau}{1+i_f\tau}.
\end{equation}
The illustrative numbers that follow are not market quotes. Take $S=1.3700$ Canadian dollars per US
dollar, a Canadian rate of 2.75\% and a US rate of 4.00\% for one year. Then $F=\cipFwd$: the currency
with the higher interest rate trades at a forward discount, here \cipPoints{} points. If a dealer quotes
\cipQuote, the US dollar is too dear for forward delivery. Borrow \cipBorrow{} Canadian dollars, buy
one million US dollars, invest them to receive \cipUSDEnd, and sell that amount forward for
\cipReceive{} Canadian dollars. The Canadian loan costs \cipOweCAD, so the trade earns \cipProfit{}
Canadian dollars with no exposure to the exchange rate. Covered parity held closely in major currencies
before 2008, but \textcite{du2018} document persistent deviations since then, which they link to the
cost of bank balance sheet space. In the terms of \cref{prop4:band}, that cost acts like a wider gap
between the rates at which an arbitrageur can borrow and lend, so the band widened.

% =====================================================================================
\section{Futures and forwards}
\label{sec4:fut}

A futures contract differs from a forward in one respect that matters for pricing. It is marked to
market at the end of every trading day, so gains and losses are paid as they happen rather than at
expiry.

\begin{definition}\label{def4:futures}
Let $0=t_0<t_1<\dots<t_n=T$ be settlement dates. A \emph{futures price} process $G_{t_k}$ with
$G_T=S_T$ defines a contract that costs nothing to enter or leave and pays the holder
$G_{t_{k+1}}-G_{t_k}$ at each $t_{k+1}$. The money market account earns the short rate $r_k$ known at
$t_k$, and its value is $B_k=\exp(\sum_{j<k}r_j(t_{j+1}-t_j))$.
\end{definition}

The payments $G_{t_{k+1}}-G_{t_k}$ are idealised variation margin. Initial margin, the interest paid on
posted collateral, funding spreads and the rules of the clearing house are left out.

The multiperiod form of \cref{thm4:ftap} says that with no arbitrage there
is a measure $Q$, equivalent to the original one, under which the gains process of every traded asset
divided by $B$ is a martingale. For an asset with no payouts that is its price divided by $B$; for an
asset with payouts it is the price divided by $B$ plus the sum of past payouts, each divided by $B$ at
its payment date. \textcite{harrison1981} prove the equivalence for finitely many states and dates, the setting of
\cref{prop4:futfwd}, building on \textcite{harrison1979}. The Vasicek and Schwartz examples below take
such a $Q$ as given in continuous time, where \textcite{harrison1981} show that pricing under it leaves
no arbitrage among admissible trading strategies. A futures position costs nothing to enter, so its gains process is just the running sum of its
settlement payments, each divided by $B$ at its payment date; that is the process to which the proof of
\cref{prop4:futfwd} applies the martingale property.

\begin{proposition}[Futures and forward prices]\label{prop4:futfwd}
Under $Q$,
\begin{equation}\label{eq4:futfwd}
  G_0=\E_Q[S_T],\qquad F_0=\frac{\E_Q[S_T/B_n]}{\E_Q[1/B_n]},\qquad
  G_0-F_0=-\frac{\Cov_Q(1/B_n,\,S_T)}{P(0,T)},
\end{equation}
where $P(0,T)=\E_Q[1/B_n]$ is the price of a zero-coupon bond. In particular $G_0=F_0$ when interest
rates are deterministic.
\end{proposition}

\begin{proof}
At $t_k$ the next settlement payment has value $\E_Q[(G_{t_{k+1}}-G_{t_k})B_k/B_{k+1}\mid\mathcal
F_{t_k}]$, and this must be zero because the contract costs nothing. The ratio $B_k/B_{k+1}=e^{-r_k(t_{k+1}-t_k)}$
is known at $t_k$ and positive, so $\E_Q[G_{t_{k+1}}\mid\mathcal F_{t_k}]=G_{t_k}$. The futures price
is a $Q$-martingale ending at $S_T$, which gives the first formula. A forward with delivery price $K$
pays $S_T-K$ at $T$ and is worth $\E_Q[(S_T-K)/B_n]$; setting this to zero gives the second. The third
follows from $\E_Q[S_T/B_n]=\E_Q[S_T]\E_Q[1/B_n]+\Cov_Q(1/B_n,S_T)$. If rates are deterministic,
$1/B_n$ is a constant and the covariance is zero.
\end{proof}

By the third formula in \cref{eq4:futfwd}, the futures price is above the forward price exactly when
$\Cov_Q(1/B_n,S_T)<0$, that is when the asset tends to end high on paths where rates have been high and
the discount factor is low. The intuition is due to \textcite{cox1981}. A long futures position gains
cash when the price rises. If rates tend to be high at the same time, the cash is reinvested at high
rates, and losses are financed at low ones. That is worth paying for. A closed form shows how large
the effect is.

\begin{proposition}[Vasicek short rate]\label{prop4:vasicek}
Under one risk-neutral measure $Q$ let $W^r$ and $W^S$ be Brownian motions with
$d\langle W^r,W^S\rangle_t=\rho\,dt$, let the short rate follow the model of \textcite{vasicek1977},
$dr_t=\kappa(\bar r-r_t)\,dt+\sigma_r\,dW^r_t$, and let an asset that pays no dividends follow
$dS_t=r_tS_t\,dt+\sigma_SS_t\,dW^S_t$. With $\Lambda(u,T)=(1-e^{-\kappa(T-u)})/\kappa$ and continuous
settlement,
\begin{equation}\label{eq4:ffratio}
  \ln\frac{G_0}{F_0}=\sigma_r^2\int_0^T\Lambda(u,T)^2\,du+\rho\,\sigma_S\sigma_r\int_0^T\Lambda(u,T)\,du .
\end{equation}
\end{proposition}

\begin{proof}
From the solution of the Ornstein-Uhlenbeck equation in Chapter~3,
$r_t=m(t)+\sigma_r\int_0^te^{-\kappa(t-u)}\,dW^r_u$ with $m$ deterministic. Integrating over
$[0,T]$ and exchanging the order of integration,
$I=\int_0^Tr_t\,dt=\bar I+\sigma_r\int_0^T\Lambda(u,T)\,dW^r_u$ with $\bar I=\int_0^Tm(t)\,dt$, which is
Gaussian with variance $v_I=\sigma_r^2\int_0^T\Lambda^2\,du$. By It\^o's lemma
$S_T=S_0\exp(I-\tfrac12\sigma_S^2T+\sigma_SW^S_T)$, and $(I,W^S_T)$ is jointly Gaussian with
$\Cov(I,\sigma_SW^S_T)=\rho\sigma_S\sigma_r\int_0^T\Lambda\,du=:\chi$. For a Gaussian $Y$,
$\E[e^Y]=e^{\E Y+\Var Y/2}$, so $G_0=\E_Q[S_T]=S_0\exp(\bar I+\tfrac12v_I+\chi)$ and
$P(0,T)=\E_Q[e^{-I}]=\exp(-\bar I+\tfrac12v_I)$. The discounted asset $e^{-I}S_T$ has mean $S_0$,
so $F_0=S_0/P(0,T)=S_0\exp(\bar I-\tfrac12v_I)$. Dividing gives \cref{eq4:ffratio}.
\end{proof}

In this model, even with $\rho=0$ the futures price is higher, because an asset whose drift is the
short rate is correlated with the discount factor through that drift. With $\kappa=0.3$,
$\sigma_r=1\%$ and $\sigma_S=20\%$, \cref{fig4:futfwd} and \cref{tab4:futfwd} show the size of the
gap in basis points of $\ln(G_0/F_0)$. The two prices agree when $\rho=\ffRhoStar$ at five years.

\begin{table}[t]
\centering
\caption{$\ln(G_0/F_0)$ in basis points from \cref{eq4:ffratio}, $\kappa=0.3$, $\sigma_r=1\%$,
$\sigma_S=20\%$.}
\label{tab4:futfwd}
\begin{tabular}{lrrr}
\toprule
maturity & $\rho=-0.5$ & $\rho=0$ & $\rho=0.5$\\
\midrule
1 year  & \ffOneNeg  & \ffOneZero  & \ffOnePos\\
5 years & \ffFiveNeg & \ffFiveZero & \ffFivePos\\
10 years& \ffTenNeg  & \ffTenZero  & \ffTenPos\\
\bottomrule
\end{tabular}
\end{table}

A simulation with \ffPaths{} paths under $Q$ (random stream 3 of the chapter seed) checks the formula.
Each path uses exact Gaussian transitions for the short rate, a trapezoidal approximation for its time
integral, and correlated Gaussian shocks for the asset, with correlation $\rho$ in each monthly step.
The monthly approximation does not reproduce the exact joint law of the integrated short rate and the
asset shock. The simulation estimates
$G_0=\E_Q[S_T]$ and $F_0=\E_Q[S_T/B_n]/\E_Q[1/B_n]$ from \cref{prop4:futfwd}, so the monthly grid
affects the approximation to the integrated short rate and its joint dependence with the asset shock;
exact short-rate transitions alone do not make the joint simulation exact. It gives \ffMCFivePos{} basis points (standard error \ffMCFivePosSE) at five years
with $\rho=0.5$, against \ffFivePos{} from the formula, and \ffMCFiveNeg{} (standard error
\ffMCFiveNegSE) against \ffFiveNeg{} with $\rho=-0.5$. Both gaps are within one standard error, which
is consistent with simulation noise, though agreement of this size does not by itself show that the
discretised rate integral adds no bias. For contracts of a few months the difference is
a basis point or two and is usually ignored. For long-dated contracts, and for interest rate futures
whose underlying is itself a rate, it is the \emph{convexity adjustment} that Chapter~6 returns to.

\begin{figure}[t]
  \centering
  \includegraphics[width=0.86\textwidth]{ch04/fig04_futfwd}
  \caption{The gap between futures and forward prices from \cref{eq4:ffratio} with $\kappa=0.3$,
  $\sigma_r=1\%$ and $\sigma_S=20\%$ in the Vasicek model, for three correlations between the asset and
  the short rate. Lines are the formula; dots are simulations with \ffPaths{} paths and monthly steps,
  with bars of two standard errors from the sample covariance of the simulated quantities (delta
  method for the log ratio).}
  \label{fig4:futfwd}
\end{figure}

% =====================================================================================
\section{Commodity term structures and the roll}
\label{sec4:roll}

A futures curve is in \emph{contango} when prices rise with maturity and in \emph{backwardation} when
they fall. Under cash and carry with a constant convenience yield, the curve is a single exponential,
up or down for every maturity. Real curves bend: a shortage pushes the front contracts up while
deferred contracts barely move, because the market expects the shortage to end. The simplest model with
that feature lets the spot price revert to a long-run level.

\begin{definition}\label{def4:schwartz}
In the one-factor model of \textcite{schwartz1997} the log spot price $X_t=\ln S_t$ is an
Ornstein-Uhlenbeck process. Under the physical measure $\mathbb P$ and the risk-neutral measure $Q$,
\begin{equation*}
  dX_t=\kappa(\alpha-X_t)\,dt+\sigma\,dW^P_t,\qquad dX_t=\kappa(\alpha^*-X_t)\,dt+\sigma\,dW^Q_t,
\end{equation*}
where the difference $\alpha-\alpha^*$ reflects a risk premium. Interest rates are deterministic.
\end{definition}

\begin{proposition}[Futures curve]\label{prop4:schwartz}
In the model of \cref{def4:schwartz}, the futures price for delivery $\tau$ years ahead is
\begin{equation}\label{eq4:schwartz}
  \ln G(X,\tau)=e^{-\kappa\tau}X+(1-e^{-\kappa\tau})\alpha^*+\frac{\sigma^2}{4\kappa}\bigl(1-e^{-2\kappa\tau}\bigr).
\end{equation}
Its slope at the short end is $\partial_\tau\ln G|_{\tau=0}=\kappa(\alpha^*-X)+\tfrac12\sigma^2$, and
$G\to\exp(\alpha^*+\sigma^2/(4\kappa))$ as $\tau\to\infty$.
\end{proposition}

\begin{proof}
By \cref{prop4:futfwd} with deterministic rates, $G=\E_Q[e^{X_{t+\tau}}\mid X_t=X]$. Under $Q$,
$X_{t+\tau}$ given $X_t=X$ is Gaussian with mean $e^{-\kappa\tau}X+(1-e^{-\kappa\tau})\alpha^*$ and
variance $\sigma^2(1-e^{-2\kappa\tau})/(2\kappa)$ (Chapter~3). The mean of a lognormal variable is the
exponential of its log mean plus half its log variance, which is \cref{eq4:schwartz}. Differentiating
at $\tau=0$ and letting $\tau\to\infty$ give the last two statements.
\end{proof}

\Cref{fig4:term} uses $\kappa=1.2$, $\sigma=0.35$ and $e^{\alpha^*}=70$. Every curve converges to
\schFinf. With the spot at 50 the curve is in contango and the one year futures price is \schOneLow;
with the spot at 100 it is in backwardation at \schOneHigh. The short end is flat when the spot is
\schSFlat. The convenience yield is not a separate input. When the spot is above \schSFlat{} the
curve slopes down because the model expects the spot to fall back, and read through
\cref{def4:convenience} the same curve shows a high implied convenience yield.

\begin{figure}[t]
  \centering
  \includegraphics[width=0.86\textwidth]{ch04/fig05_term}
  \caption{Futures curves from \cref{eq4:schwartz} for five spot prices. Low spot prices give contango
  and high ones backwardation, and all curves approach the same long-maturity level.}
  \label{fig4:term}
\end{figure}

\paragraph{Rolling a position} An investor who wants commodity exposure for years cannot hold one
contract, because each expires. The usual strategy holds a near contract, sells it shortly before
expiry and buys the next one. Write $b_t(T)=\ln G_t(T)-\ln S_t$ for the \emph{basis} of the contract
expiring at $T$.

\begin{proposition}[Spot and roll]\label{prop4:roll}
Hold the contract expiring at $T_k$ from $t_k$ to $t_k+h$, then move to the next contract, for
$k=0,\dots,m-1$, with $t_{k+1}=t_k+h$. The sum of the log changes in the futures prices held splits
exactly as
\begin{equation}\label{eq4:roll}
  \sum_{k}\ln\frac{G_{t_k+h}(T_k)}{G_{t_k}(T_k)}
  =\underbrace{\ln\frac{S_{t_m}}{S_{t_0}}}_{\text{spot}}
  +\underbrace{\sum_k\bigl[b_{t_k+h}(T_k)-b_{t_k}(T_k)\bigr]}_{\text{roll}} .
\end{equation}
\end{proposition}

\begin{proof}
For each $k$, $\ln G_{t_k+h}(T_k)-\ln G_{t_k}(T_k)=(\ln S_{t_k+h}-\ln S_{t_k})+(b_{t_k+h}(T_k)-b_{t_k}(T_k))$
by the definition of $b$. Summing, the spot terms telescope.
\end{proof}

The left side of \cref{eq4:roll} is a statistic of futures prices. Money invested in the strategy
grows at a slightly different rate, and the difference can be computed exactly. Consider an investor who keeps the
whole value $V_k$ of the position in bills earning a constant rate $r$ and holds $V_k/G_{t_k}(T_k)$
contracts of unit size from $t_k$ to $t_k+h$, so that the futures exposure equals current wealth, resetting the number of contracts at each step. Variation margin is
paid out of, or into, the bills, and fees, funding spreads and initial margin are ignored. Write
$g_k=\ln\bigl(G_{t_k+h}(T_k)/G_{t_k}(T_k)\bigr)$. After one step the bills are worth $V_ke^{rh}$ and
the contracts have paid $V_k(e^{g_k}-1)$, so
\begin{equation}\label{eq4:wealth}
  V_{k+1}=V_k\bigl(e^{rh}+e^{g_k}-1\bigr),\qquad
  \ln\frac{V_m}{V_0}=\sum_k\bigl(rh+g_k\bigr)+\sum_k\ln\bigl(1-(1-e^{-g_k})(1-e^{-rh})\bigr),
\end{equation}
where the logarithm exists because the wealth factor $e^{rh}+e^{g_k}-1$ is positive whenever $r\ge0$,
since $e^{g_k}>0$; only with a negative rate could a near-total fall in the futures price in one
step make it nonpositive. The second form uses
$e^{rh}+e^{g}-1=e^{rh+g}\bigl(1-(1-e^{-g})(1-e^{-rh})\bigr)$. The log
growth of the collateralised position is therefore the bill return plus the sum of log futures
changes, plus a correction whose terms are close to $rh\,(e^{-g_k}-1)$, the product of the bill
return and the futures move. Treating the sum in \cref{eq4:roll} plus $r$ as the strategy's return is
an approximation that drops this correction. On the simulated paths below, with $r=\rollRate\%$ and
monthly steps, the correction is \rollGapMedBp{} basis points a year at the median across paths and
at most \rollGapMaxBp{} basis points a year.

The identity in \cref{eq4:roll} is exact and holds in any model; what a model adds is the size of the
roll. In backwardation the held contract sits below spot and, if the curve keeps its shape, rises
towards it as expiry nears, so the roll tends to be positive. In contango it tends to be negative.

\begin{proposition}[Expected roll]\label{prop4:expectedroll}
In the model of \cref{def4:schwartz}, with $X$ stationary under the physical measure $\mathbb P$ and contracts
held from $\tau$ to $\tau-h$ years before expiry, the expected roll per step and the expected log
futures change per step are both
\begin{equation}\label{eq4:expectedroll}
  \bigl(e^{-\kappa(\tau-h)}-e^{-\kappa\tau}\bigr)(\alpha-\alpha^*)
  -\frac{\sigma^2}{4\kappa}\bigl(e^{-2\kappa(\tau-h)}-e^{-2\kappa\tau}\bigr).
\end{equation}
\end{proposition}

\begin{proof}
By \cref{eq4:schwartz}, $b(X,\tau)=\ln G(X,\tau)-X$ is affine in $X$ with slope $e^{-\kappa\tau}-1$.
Under stationarity $\E_{\mathbb P}[X_t]=\alpha$ at every date, so the expected roll is
$b(\alpha,\tau-h)-b(\alpha,\tau)$, which simplifies to \cref{eq4:expectedroll}. The expected spot term
$\E_{\mathbb P}[X_{t+h}-X_t]$ is zero, so the expected log futures change is the same. Both statements are about
the model's decomposition of log prices under $\mathbb P$; the expected growth of wealth adds $r$ and the
correction in \cref{eq4:wealth}.
\end{proof}

With $\alpha-\alpha^*=0.05$, monthly rolls from two months to one month before expiry, and the
parameters above, \cref{eq4:expectedroll} gives \rollTheory\% a year, a small difference of two larger terms: the premium term contributes
\rollPremTerm\% a year and the convexity term \rollConvTerm\%. The first term is the risk
premium; the second is the convexity cost of measuring returns in logs, and it would vanish in
arithmetic returns. A simulation of \rollPaths{} independent forty year paths under $\mathbb P$, each started
from the stationary law and moved by exact Ornstein-Uhlenbeck transitions (random stream 4 of the
chapter seed), gives an average roll of \rollMCMean\% a year (standard error \rollMCSE). Their
average log wealth growth from \cref{eq4:wealth}, bills included, is \rollWealthMean\% a year, against
\rollApproxMean\% from the sum of log futures changes plus $r$. For a single investor the spread across paths matters more. The spot term telescopes to the change
in the log spot price over the whole period, which stationarity keeps small, while the roll adds a
term every month that depends on the basis at that date, so a path that spends years in contango or
in backwardation accumulates it. Across paths the forty year average roll ranges from \rollMCLo\% to
\rollMCHi\% a year (5th to 95th percentile), while the average
spot return ranges only from \rollMCSpotLo\% to \rollMCSpotHi\%. The cross-path standard deviation of the average roll, \rollMCSd{} points,
measures how far investors' outcomes differ; it is not a standard error. In this stationary model, then, the spread of long-run outcomes for a
rolled position comes mostly from the roll.

\Cref{fig4:roll} shows the path whose average roll is the median of the \rollPaths. Its spot price
fell by \rollSpotAnnAbs\% a year in logs, a loss of \rollSpotFortyFall\% over the forty years; the curve was in contango in \rollContango\% of
months, and the roll added \rollRollAnn\% a year, close to the model's \rollTheory\%. The
collateralised position grew at \rollWealthAnn\% a year in logs, against \rollTotalAnn\% from the sum
of log futures changes plus $r$. Panel (b) illustrates, for this particular path, how changes in the basis feed the measured roll: on this path the roll
in a month has a correlation of \rollBasisCorr{} with the basis at entry.

Whether rolled commodity positions earn a premium in practice is an empirical question, separate
from the model. \textcite{gorton2006} find long-run returns on a collateralised index comparable to
equities, and \textcite{erb2006} argue that differences in roll returns explain much of the difference
in long-run returns across commodities. The older theory of normal backwardation
\parencite{keynes1930} explains a positive premium as payment to speculators who take risk off
producers' hands, and $\alpha>\alpha^*$ is one way to put that compensation into the model, with no claim that it measures the premium in actual markets.

\begin{figure}[t]
  \centering
  \includegraphics[width=\textwidth]{ch04/fig06_roll}
  \caption{(a) The decomposition of \cref{eq4:roll} along the simulated forty year path whose average
  roll is the median of \rollPaths{} paths, with monthly rolls. (b) Monthly roll against the basis
  $\ln(G/S)$ at entry on that path: a negative basis at entry (backwardation) goes with a positive
  roll.}
  \label{fig4:roll}
\end{figure}

% =====================================================================================
\section{The binomial model}
\label{sec4:binom}

\begin{definition}\label{def4:binomial}
In the \emph{binomial model} with $n$ steps a stock starts at $S_0$ and at each step is multiplied by
$u$ or by $d$, with $0<d<u$. A riskless account grows by the factor $R>0$ per step. A portfolio holds
$\Delta_k$ shares and $\beta_k$ units of currency in the account from step $k$ to step $k+1$, chosen at
step $k$ with knowledge of the path so far; at step $k+1$, before rebalancing, it is worth
$\Delta_kS_{k+1}+\beta_kR$. It is \emph{self-financing} if
$\Delta_kS_{k+1}+\beta_kR=\Delta_{k+1}S_{k+1}+\beta_{k+1}$ for $k=0,\dots,n-2$, so that rebalancing at
step $k+1$ neither adds nor removes money.
\end{definition}

\begin{proposition}[One step]\label{prop4:onestep}
The one-step market has no arbitrage if and only if $d<R<u$. In that case a claim paying $V_u$ after
an up move and $V_d$ after a down move is replicated by
\begin{equation}\label{eq4:onestep}
  \Delta=\frac{V_u-V_d}{S(u-d)},\qquad \beta=\frac{uV_d-dV_u}{R(u-d)},
\end{equation}
and its price is $V=\Delta S+\beta=\bigl(qV_u+(1-q)V_d\bigr)/R$ with $q=(R-d)/(u-d)$.
\end{proposition}

\begin{proof}
The market is complete with two assets and two states. By \cref{thm4:ftap} there is no arbitrage if and
only if the unique solution of $\psi_u+\psi_d=1/R$ and $\psi_uSu+\psi_dSd=S$ is positive. The solution
is $\psi_u=q/R$ and $\psi_d=(1-q)/R$, positive exactly when $0<q<1$, that is $d<R<u$. Solving
$\Delta Su+\beta R=V_u$ and $\Delta Sd+\beta R=V_d$ gives \cref{eq4:onestep}, and the price is
$\psi_uV_u+\psi_dV_d$ by \cref{cor4:complete}.
\end{proof}

Take $S=100$, $u=1.2$, $d=0.9$ and $R=1.05$, so $q=\opQ$. A call struck at 100 pays 20 or 0. The
replicating portfolio holds $\Delta=\opDelta$ shares and $\beta=\opBond$ in the account, which is a loan, and costs \opPrice, the same
as $(q\cdot20)/R$. The probability that the stock actually rises plays no part: two investors who
disagree about it still agree on the price of the call, because both can build it from the stock and
the account.

\begin{theorem}[Risk-neutral pricing on a tree]\label{thm4:tree}
Suppose $d<R<u$ and let $Q$ make the moves independent with up probability $q=(R-d)/(u-d)$. A claim
paying $H$ at step $n$, where $H$ may depend on the whole path, is replicated by a self-financing
portfolio whose value at step $k$ is $V_k=R^{-(n-k)}\E_Q[H\mid\mathcal F_k]$. If $H=g(S_n)$,
\begin{equation}\label{eq4:binsum}
  V_0=R^{-n}\sum_{j=0}^n\binom nj q^j(1-q)^{n-j}\,g\bigl(S_0u^jd^{n-j}\bigr).
\end{equation}
\end{theorem}

\begin{proof}
Define $V_n=H$ and $V_k=R^{-1}\E_Q[V_{k+1}\mid\mathcal F_k]$ for $k<n$. At a node at step $k$ the next
value takes two values, $V_{k+1}^u$ and $V_{k+1}^d$, and $V_k=(qV_{k+1}^u+(1-q)V_{k+1}^d)/R$ is the price
of a one-step claim. \Cref{prop4:onestep} gives $(\Delta_k,\beta_k)$ with $\Delta_kS_k+\beta_k=V_k$ and
$\Delta_kS_{k+1}+\beta_kR=V_{k+1}$ in both successor nodes. The second equation is the self-financing
condition with $\Delta_{k+1}S_{k+1}+\beta_{k+1}=V_{k+1}$. So at every node the old portfolio pays for the new one, and the strategy carries $V_0$
to $H$ without further cash. By the tower property $V_k=R^{-(n-k)}\E_Q[H\mid\mathcal F_k]$, and under $Q$ the number of
up moves is binomial with parameters $n$ and $q$, which gives \cref{eq4:binsum}.
\end{proof}

% =====================================================================================
\section{From trees to the Black-Scholes price}
\label{sec4:limit}

To approximate a stock with volatility $\sigma$ over $[0,T]$, divide time into $n$ steps of length
$h=T/n$ and set
\begin{equation}\label{eq4:crr}
  u=e^{\sigma\sqrt h},\qquad d=e^{-\sigma\sqrt h},\qquad R=e^{rh},\qquad q_n=\frac{e^{rh}-e^{-\sigma\sqrt h}}{e^{\sigma\sqrt h}-e^{-\sigma\sqrt h}} .
\end{equation}
These are the parameters of \textcite{cox1979}. For $h$ small, $d<R<u$, so \cref{thm4:tree} applies.

\begin{theorem}[Convergence to Black-Scholes]\label{thm4:crr}
Let $S_0,K,T,\sigma>0$, let the rate $r$ be constant and let the stock pay no dividends. For each $n$
let $P_n$ and $C_n$ be the prices of a European put and call with strike $K$ and expiry $T$ on the tree
defined by \eqref{eq4:crr}. As $n\to\infty$, $P_n$ and $C_n$ converge to the Black-Scholes prices
\begin{equation}\label{eq4:bs}
  C=S_0\Phi(d_1)-Ke^{-rT}\Phi(d_2),\qquad P=Ke^{-rT}\Phi(-d_2)-S_0\Phi(-d_1),
\end{equation}
where $d_{1,2}=\bigl(\ln(S_0/K)+(r\pm\tfrac12\sigma^2)T\bigr)/(\sigma\sqrt T)$.
\end{theorem}

\begin{proof}
Under $Q$, $\ln(S_n/S_0)=\sigma\sqrt h\sum_{k=1}^n\xi_k$ with independent $\xi_k=\pm1$ and
$Q(\xi_k=1)=q_n$. Write $\varepsilon_h=\sigma\sqrt h$ for the size of one log step. Expanding numerator and denominator of $q_n$,
\begin{equation*}
  q_n=\frac{\varepsilon_h+rh-\tfrac12\varepsilon_h^2+O(h^{3/2})}{2\varepsilon_h+O(h^{3/2})}=\frac12+\frac{(r-\tfrac12\sigma^2)\sqrt h}{2\sigma}+O(h),
\end{equation*}
so $\E_Q[\xi_k]=2q_n-1=(r-\tfrac12\sigma^2)\sqrt h/\sigma+O(h)$. Hence the mean of the log return is
$n\,\varepsilon_h\,(2q_n-1)=(r-\tfrac12\sigma^2)T+O(\sqrt h)$, and its variance is $n\,\varepsilon_h^2\bigl(1-(2q_n-1)^2\bigr)=\sigma^2T+O(h)$.
The summands $\varepsilon_h\xi_k$ form a triangular array of independent variables bounded by $\varepsilon_h\to0$, so the
Lindeberg condition holds trivially, and the Lindeberg-Feller central limit theorem gives
$\ln(S_n/S_0)\Rightarrow\N\bigl((r-\tfrac12\sigma^2)T,\sigma^2T\bigr)$. The put payoff
$x\mapsto(K-S_0e^x)^+$ is bounded and continuous, so $P_n=R^{-n}\E_Q[(K-S_n)^+]$ converges to
$e^{-rT}\E[(K-S_T)^+]$ with $\ln S_T$ normal as above, which evaluates to $P$ in \cref{eq4:bs} (the
calculation is the lognormal integral of Chapter~1). Finally put-call parity holds exactly on every
tree, $C_n-P_n=S_0-KR^{-n}=S_0-Ke^{-rT}$, because the difference of payoffs is $S_n-K$; the same identity
holds for $C$ and $P$, so $C_n\to C$.

The call can also be handled directly, which matters for payoffs without a parity partner.
Convergence in distribution gives convergence of expectations for an unbounded payoff once the payoffs
are uniformly integrable. Here $\E_Q[S_n^2]=S_0^2\bigl(q_nu^2+(1-q_n)d^2\bigr)^n$ and
$q_nu^2+(1-q_n)d^2=R(u+d)-ud=2e^{rh}\cosh(\sigma\sqrt h)-1=1+(2r+\sigma^2)h+O(h^2)$, so
$\E_Q[S_n^2]\to S_0^2e^{(2r+\sigma^2)T}$ is bounded in $n$. Since $0\le(S_n-K)^+\le S_n$, the call
payoffs are bounded in $L^2(Q)$ and hence uniformly integrable, and $C_n$ converges to the discounted
lognormal expectation.
\end{proof}

Chapter~5 derives \cref{eq4:bs} directly; here it appears as a limit. \Cref{fig4:conv} shows how the
limit is reached for a one year call with $S_0=100$, $K=110$, $r=5\%$ and $\sigma=20\%$, whose
Black-Scholes price is \bsCall. The tree prices are \crrFifty{} at $n=50$, \crrHundred{} at $n=100$,
\crrHundredOne{} at $n=101$ and \crrFiveHundred{} at $n=500$. The error does not shrink smoothly. It
oscillates because the strike falls at a different place relative to the terminal nodes as $n$
changes. For a European call or put on this tree, \textcite{diener2004} show that $n$ times the error
stays bounded but does not converge, with a pattern that depends on the strike; here it moves between
\crrNErrMin{} and \crrNErrMax{} for $100\le n\le500$. Fitted to the largest absolute error in eight geometric blocks of $n$ from 20 to 500, the envelope
falls with a slope of \crrSlope{} on a log scale, close to the theoretical $-1$, although the fitted slope over this finite range is not an estimate
of the asymptotic order.

\begin{figure}[t]
  \centering
  \includegraphics[width=\textwidth]{ch04/fig08_convergence}
  \caption{(a) Tree prices of a one year call, $S_0=100$, $K=110$, $r=5\%$, $\sigma=20\%$, against the
  number of steps. (b) Absolute errors $|C_n-C|$ against the Black-Scholes price, on log scales. The
  Cox-Ross-Rubinstein error oscillates inside an envelope of order $1/n$ because the strike sits at a
  different place between terminal nodes as $n$ changes; the Leisen-Reimer tree, defined for odd $n$
  and centred on the strike, converges like $1/n^2$ here.}
  \label{fig4:conv}
\end{figure}

The oscillation can be removed. \textcite{leisen1996} take an odd number of steps, so that the middle
of the terminal nodes falls on the strike, and choose $q$ and the up and down factors so that the
binomial tail probabilities match $\Phi(d_1)$ and $\Phi(d_2)$ through the normal approximation of
\textcite{peizer1968}. For this construction and payoff they find the error falling like $n^{-2}$;
here the fitted slope for odd $n\ge31$ is \lrSlope, and at $n=101$ the price is \lrHundredOne, an error
of \lrErrHundredOne{} millionths. Back on the plain Cox-Ross-Rubinstein tree, at the root of a 1,000 step tree the hedge ratio
of \cref{eq4:onestep} for the same call is \treeDelta, equal to the Black-Scholes delta $\Phi(d_1)=\bsDelta$ to four
places.

The listing below prices a one year at-the-money put ($S_0=K=100$, $r=5\%$, $\sigma=20\%$, 1,000
steps) by backward induction, as a European option and, with one extra line, as an American one.
\lstinputlisting[language=Python]{code/ch04/snippets/crr.py}
It prints \snipEu{} and \snipAm. The European value differs from the Black-Scholes put, \bsPut, in the
third decimal; the American value is the subject of the next section.

% =====================================================================================
\section{American options}
\label{sec4:american}

An American option can be exercised at any step. Its holder chooses a \emph{stopping time}: a rule
$\tau$ taking values in $\{0,\dots,n\}$ such that whether $\tau=k$ is decided by the path up to step
$k$. For an option with exercise value $g(S_k)$, define the \emph{Snell envelope} by backward
induction:
\begin{equation}\label{eq4:snell}
  V_n=g(S_n),\qquad V_k=\max\bigl\{g(S_k),\;R^{-1}\E_Q[V_{k+1}\mid\mathcal F_k]\bigr\}.
\end{equation}

\begin{theorem}[Optimal exercise]\label{thm4:snell}
$V_0=\max_\tau\E_Q[R^{-\tau}g(S_\tau)]$, the maximum being over stopping times, and it is attained by
$\tau^*=\min\{k:V_k=g(S_k)\}$.
\end{theorem}

\begin{proof}
Let $Y_k=R^{-k}V_k$. By \cref{eq4:snell}, $Y_k\ge\E_Q[Y_{k+1}\mid\mathcal F_k]$, so $Y$ is a
supermartingale, and $Y_k\ge R^{-k}g(S_k)$. For a stopping time $\tau\le n$,
\begin{equation*}
  Y_{\tau}=Y_0+\sum_{k=0}^{n-1}\mathbf 1\{\tau>k\}(Y_{k+1}-Y_k),
\end{equation*}
and each term has $\E_Q[\mathbf 1\{\tau>k\}(Y_{k+1}-Y_k)]\le0$ because $\{\tau>k\}$ is known at step
$k$. So $\E_Q[R^{-\tau}g(S_\tau)]\le\E_Q[Y_\tau]\le Y_0=V_0$. For $\tau^*$, on $\{\tau^*>k\}$ the
maximum in \cref{eq4:snell} is the continuation value, so $Y_k=\E_Q[Y_{k+1}\mid\mathcal F_k]$ there
and every term above has expectation zero. Hence $\E_Q[Y_{\tau^*}]=V_0$, and $Y_{\tau^*}=R^{-\tau^*}g(S_{\tau^*})$
by the definition of $\tau^*$.
\end{proof}

The seller can hedge an American option too. At each step the seller holds the portfolio that
replicates next step's option value $V_{k+1}$. It costs the continuation value
$R^{-1}\E_Q[V_{k+1}\mid\mathcal F_k]\le V_k$, so any excess of $V_k$ over that can be withdrawn, and at
the next step the portfolio is worth $V_{k+1}\ge g(S_{k+1})$, enough to pay the holder if the option is
exercised there. No cash is ever added. The two-period put of \cref{fig4:tree} shows the exercise rule at work.
With $K=105$ and the tree of the example after \cref{prop4:onestep}, the
put is in the money at the down node, $S=90$, where exercise pays \tpExDown{} and waiting is worth
\tpContDown. The holder exercises, and the put is worth \tpAmerican{} today against \tpEuropean{} for
the European put, an early exercise premium of \tpPremium.

\begin{figure}[t]
  \centering
  \includegraphics[width=0.8\textwidth]{ch04/fig07_tree}
  \caption{A two-period American put with $S_0=100$, $u=1.2$, $d=0.9$, $R=1.05$ and $K=105$. Each node
  shows the stock price and the option value from \cref{eq4:snell}; exercise is optimal at the red node.}
  \label{fig4:tree}
\end{figure}

\begin{proposition}[No early exercise of a call]\label{prop4:merton}
On a tree with $R>1$ and a stock that pays no dividends, an American call with a strictly positive
strike has the same value as the
European call, and early exercise is never strictly optimal.
\end{proposition}

\begin{proof}
Let $c_k$ be the European call value at step $k$. Since $x\mapsto(x-K)^+$ is convex and
$\E_Q[S_n\mid\mathcal F_k]=R^{n-k}S_k$, Jensen's inequality gives
$c_k\ge R^{-(n-k)}(R^{n-k}S_k-K)^+\ge S_k-KR^{-(n-k)}$. Since also $c_k\ge0$, this gives
$c_k\ge(S_k-K)^+$ at every node. Now argue backwards: if $V_{k+1}=c_{k+1}$, then the continuation value
in \cref{eq4:snell} is $c_k$, which is at least the exercise value, so $V_k=c_k$. That proves the
equality of the two values. For the strict statement, take $k<n$ and a node with $S_k>K$: as $R>1$ and
$K>0$, $S_k-KR^{-(n-k)}>S_k-K$, so $c_k>(S_k-K)^+$ and exercising there gives up value. At a node with
$S_k\le K$ exercise pays nothing, so early exercise is never strictly optimal.
\end{proof}

This is the tree version of a result of \textcite{merton1973}. A dividend breaks the argument, because
$\E_Q[S_n\mid\mathcal F_k]$ is then less than $R^{n-k}S_k$. With a continuous dividend yield of 4\%,
which enters the tree through $q=(e^{(r-\delta)h}-d)/(u-d)$ while values are discounted at $r$, the one
year at-the-money call ($r=5\%$, $\sigma=20\%$) is worth \amCallDiv{} as an American option and \euCallDiv{} as a European one on a
1,000 step tree. Without the dividend the two trees agree to within $10^{-10}$, both at \amCallNoDiv.

For puts early exercise is valuable even without dividends: receiving $K$ early earns interest. The
next result shows that the region where exercise is optimal is always below a boundary.

\begin{proposition}[Exercise boundary of a put]\label{prop4:boundary}
On any recombining tree with $d<R<u$ and no dividends, let $V_{k,j}$ be the American put value at the
node with $j$ up moves after $k$ steps. Then $V_{k,j}+S_{k,j}$ is nondecreasing in $j$. Consequently,
at each step at which exercise is optimal at some node there is a level $b_k$, one of the node prices of that step, such that the put's exercise region at positive payoff consists
exactly of the nodes with $S_{k,j}\le b_k$ (and $b_k<K$).
\end{proposition}

\begin{proof}
Let $w_{k,j}=V_{k,j}+S_{k,j}$. At step $n$, $w_{n,j}=\max(K,S_{n,j})$, which is nondecreasing in $j$.
Suppose the claim holds at step $k+1$. Because the stock pays no dividends,
$S_{k,j}=R^{-1}(qS_{k+1,j+1}+(1-q)S_{k+1,j})$, and \cref{eq4:snell} becomes
\begin{equation*}
  w_{k,j}=\max\bigl\{K,\;R^{-1}\bigl(q\,w_{k+1,j+1}+(1-q)\,w_{k+1,j}\bigr)\bigr\},
\end{equation*}
which is nondecreasing in $j$ because both terms inside the maximum are. At a node in the money,
$S_{k,j}<K$, immediate exercise is optimal if and only if $w_{k,j}=K$. Since $w_{k,j}\ge K$ and is
nondecreasing, $\{j:w_{k,j}=K\}$ is an initial segment of the indices; together with $S_{k,j}<K$ it gives the exercise set. If the set is
empty there is nothing to prove, and we set $b_k=0$ by convention; otherwise its largest index is
$j_k$ and $b_k=S_{k,j_k}$.
\end{proof}

\Cref{fig4:boundary} plots the boundary for a put with $K=100$, $r=5\%$ and $\sigma=20\%$. The levels
$b_k$ are node prices. Because a finite tree admits exercise only at its nodes, successive estimates of
the exercise threshold alternate between nearby stock prices. The plotted lower monotone envelope in
time to expiry is a picture of that discrete boundary; it is not computed separately as a
continuous-time exercise boundary, and the levels quoted below are envelope estimates unless said otherwise. A one year tree has only
its root one year before expiry, so the one year value comes from a two year, 4,000 step tree with the same
step of 0.0005 years, which has nodes there. One year before expiry the lower-envelope estimate of the tree's exercise threshold is
\bndStart, six months before it is \bndHalf, and close to expiry the boundary rises to the strike.
On a ten year, 2,000 step tree the lowest node first falls into the exercise region
\bndTenFirstT{} years before expiry, at \bndTenFirst; at earlier steps no node is low enough, so that
date reflects the grid and not the option. From there the raw levels alternate between neighbouring
nodes: one step later, \bndTenSecondT{} years before expiry, the level is already the lower node,
\bndTenSecond. The envelope, the running minimum taken outward from expiry, first reaches \bndTenFirstEnv{}
\bndTenEnvT{} years before expiry, the step nearest expiry at which that node is the level, and stays
there out to ten years, approaching the boundary of the perpetual put, which \textcite{merton1973} found in closed form. These
are properties of this model, with constant rate and volatility and no dividends. In the continuous model a perpetual put does not depend on time, so its value solves
$\tfrac12\sigma^2S^2V''+rSV'-rV=0$ where it is not exercised. The solutions are $S$ and $S^{-\gamma}$
with $\gamma=2r/\sigma^2$; boundedness excludes the first, and the conditions $V(b)=K-b$ and
$V'(b)=-1$ at the boundary give $b=\gamma K/(1+\gamma)=\perpBoundary$. The second condition, smooth
fit, is taken here as given; \textcite{shreve2004} treats the perpetual put in full.

\begin{figure}[t]
  \centering
  \includegraphics[width=0.86\textwidth]{ch04/fig09_boundary}
  \caption{The level below which an American put should be exercised, against time to expiry, for a one
  year put on a 2,000 step tree and a ten year put on 2,000 steps; the one year level quoted in the
  text comes from a two year, 4,000 step tree with the same step of 0.0005 years. Node prices form a grid, so the boundary is shown as the
  lower monotone envelope of the levels of \cref{prop4:boundary}; the envelope hides the step-to-step
  alternation between neighbouring nodes. The dashed line is the perpetual put boundary.}
  \label{fig4:boundary}
\end{figure}

The one year American put is worth \amPutThousand{} on a 1,000 step tree and \amPutTwoThousand{} on
2,000 steps, against \euPutTwoThousand{} for the European put on the same 2,000 step tree, so the right to exercise
early is worth \amPremiumTree. For comparison, the Black-Scholes European value is \bsPut; subtracting
it from the tree American price gives \amPremium, a mixed-method figure and not a same-tree premium. There is no generally used elementary closed form for the finite-maturity American
put under these assumptions; trees, together with the finite difference methods of Chapter~5, are a
standard way to price it.

% =====================================================================================
\section{Case study: do forward rates forecast?}
\label{sec4:case}

\Cref{prop4:fra} shows that anyone can lock in a forward rate with two bonds. Whether it is also a
good forecast depends on the term premium, the extra return that investors demand for holding longer
bonds. The test of \textcite{fama1987} regresses the realised change in the spot rate on the
forward-spot spread:
\begin{equation}\label{eq4:fb}
  y_{t+12}-y_t=a+b\,\bigl(F_t(1,2)-y_t\bigr)+\varepsilon_{t+12},
\end{equation}
where $y_t$ is the one year bill rate at the end of month $t$, a simple rate, and $F_t(1,2)$ is the
simple forward rate, read from that month's curve, for borrowing from one year ahead to two years
ahead. The pure expectations hypothesis gives $a=0$ and $b=1$. A constant premium keeps $b=1$ and makes
$a$ equal to minus the premium. Both statements hold to first order, because the hypothesis is stated
for expected returns while the regression uses simple rates. A slope of zero means the spread carries
no information about the future change.

The data are month-end Government of Canada yields from \fbStart{} to \fbEnd, \caseMonths{} months,
from the Bank of Canada's Valet service \parencite{bocvalet}. The response is the change in the one
year bill rate over the next twelve months and the explanatory variable is the one year forward rate
one year ahead less the current one year bill rate, both at month ends. Each month the six month and one year bills and the
two and three year benchmark bonds are bootstrapped as in \cref{sec4:curve}, and the forward rates are
read off the curve on the same simple basis as the bill yield. Each month is matched with the bill
rate twelve months later, so the last twelve months have no realised outcome and drop out, leaving
\fbOneN{} observations for the regression; the last month-end is \fbOneLastDate, observed on \gocDate. The regression
errors overlap by eleven months, so standard errors use the estimator of \textcite{newey1987} with
twelve lags.

The slope is \fbOneSlope{} with a standard error of \fbOneSE{} on \fbOneN{} observations, a 95\%
interval of \fbOneLo{} to \fbOneHi. Zero is rejected ($t=\fbOneTZero$) and one is not, though narrowly
($t=\fbOneTOne$). Failing to reject one is weaker than evidence that the slope is one: the interval
also covers slopes near a quarter. The forward-spot spread explains \fbOneRSq{} of the variance of the
change. Two years ahead, regressing $y_{t+24}-y_t$ on $F_t(2,3)-y_t$ with 24 Newey-West lags, the slope
is \fbTwoSlope{} (standard error \fbTwoSE, interval \fbTwoLo{} to \fbTwoHi) on \fbTwoN{} observations.
That regression loses two years at the end and its errors overlap by 23 months, so it is less precise. On average the forward rate
exceeded the rate it targets by \fwdErrMean{} basis points, with a standard error of \fwdErrSE, so the average error is not distinguishable from zero
($t\approx\fwdErrT$), and it
was too high in \fwdAbove\% of months. \Cref{fig4:case} shows why the evidence is weak in both
directions. Most of the variation in the change comes from a few episodes, the large rate changes around 2008,
2020 and 2022, which the estimated forward rates did not predict.

\begin{figure}[t]
  \centering
  \includegraphics[width=\textwidth]{ch04/fig10_case}
  \caption{(a) The one year Government of Canada bill rate at month ends from \fbStart{} to
  \fbEnd, and the forward rate observed twelve months earlier, plotted at the date it targets, so its
  line runs from \fbFwdStart{} to \fbFwdEnd. The forward rate is a market rate that includes
  any term premium, so it need not be an unbiased forecast. (b) The regression \cref{eq4:fb}: the fitted slope
  is \fbOneSlope, below the slope of one implied by a constant term premium, though one lies inside its 95\%
  interval.}
  \label{fig4:case}
\end{figure}

Three cautions apply. Twenty-five years of overlapping annual changes contain about 25 non-overlapping annual windows,
which need not be 25 statistically independent observations, and conventional Newey-West inference can be unreliable when observations overlap this heavily, so the $t$ statistic against one should not be read precisely. Benchmark bonds are not par bonds
of exact maturity, so the forward rates carry measurement error, which, if it is unrelated to the true spread and to the
later rate change, biases the slope towards zero.
And as a forecast in the plain sense the forward rate did worse than assuming no change: its mean
absolute error a year ahead was \fwdErrMAE{} basis points against \naiveMAE{} for the no-change
forecast. The slope is positive but imprecise over this sample. What it leaves
unexplained is consistent with a term premium that varies over time, a premium the regression itself
does not measure. \textcite{fama1987}
found the slope rising with the horizon in United States data; here the one and two year slopes are
almost equal (\fbOneSlope{} and \fbTwoSlope), so this sample does not show that pattern.

% =====================================================================================
\section{Summary}
\label{sec4:summary}

A forward rate, a forward price and an option price on a tree all come from the same argument: build
the payoff from things that already trade, and charge what they cost. In a one-period market that
argument works exactly when there are positive state prices. When the market is incomplete, a claim
that can be replicated still has one price, and a claim
that cannot be replicated has an open interval of arbitrage-free prices between its sub- and
super-replication bounds.

Cash and carry fixes the forward price of a financial asset at the spot price grown at the interest
rate less the income yield. Frictions turn the price into a band. For currencies the frictionless
relation is covered interest parity, and transaction costs and funding spreads widen it into a band of
the same kind. Commodities add storage costs and a convenience yield that only holders of inventory enjoy, so only the cash-and-carry side of the
arbitrage works. Futures and forwards differ when interest rates are random, by an amount that
depends on the covariance between the asset and the discount factor and, in the Vasicek example, grows
with maturity. For a rolled futures position the sum of log futures changes is the spot return plus
the roll, and it approximates the log return of a collateralised position less the bill rate. In the
stationary commodity model of \cref{sec4:roll}, the roll accounts for most of the spread of long-run
outcomes.

On a binomial tree every claim is replicated by a self-financing portfolio and priced as a
discounted risk-neutral expectation. As the steps shrink, the price converges to the Black-Scholes
formula; for a call or put on the Cox-Ross-Rubinstein tree the error oscillates inside an envelope of
order $1/n$ unless the tree is tuned. American options are priced by the Snell envelope, and calls on
stocks without dividends are never exercised early. For the constant-parameter model without dividends, the put's exercise threshold falls as the
remaining maturity lengthens. In the Canadian data the forecasting
slope of forward rates is positive but imprecisely estimated.

% =====================================================================================
\section{Exercises}
\label{sec4:exercises}

\begin{enumerate}[label=\textbf{4.\arabic*}, leftmargin=*, ref=4.\arabic*]
  \item\label{ex4:one} A one year bill yields 3.00\% on a simple basis and a two year bond with an
  annual coupon has a par yield of 3.50\%. Find $D(1)$, $D(2)$ and the forward rate $F(1,2)$.
  \item\label{ex4:two} An index trades at 50 with a dividend yield of 2\% and the interest rate is 5\%,
  both continuous. Find the nine month forward price.
  \item\label{ex4:three} A long forward on that index has a delivery price of 52 and four months to run.
  The index is at 55. What is the contract worth?
  \item\label{ex4:four} The spot rate is 1.35 Canadian dollars per US dollar, and six month simple rates
  are 3\% in Canada and 5\% in the United States. Find the parity forward rate. A dealer quotes 1.3400;
  describe the arbitrage on one million US dollars and find its profit in Canadian dollars.
  \item\label{ex4:five} A stock at 100 will be at 110 or 90 next period and the interest rate is zero.
  Find the state prices and the prices of a call struck at 100 and a put struck at 105.
  \item\label{ex4:six} A commodity trades at 70 spot and 66 for delivery in one year. The interest rate is
  4\% and storage costs 2\%, both continuous. Find the implied convenience yield.
  \item\label{ex4:seven} On a one-step tree with $S=50$, $u=1.1$, $d=0.9$ and $R=1.02$, price a put struck
  at 50 and find the number of shares in its replicating portfolio.
  \item\label{ex4:eight} Extend Exercise~\ref{ex4:seven} to two steps, with the same $u$, $d$ and $R$
  at each step, and strike 52. Price the put as European and as American, with exercise allowed at
  steps 0, 1 and 2. At which node is it exercised?
  \item\label{ex4:nine} A trader buys the front futures contract at 72 with the spot at 70 and sells it a
  month later at 74 with the spot at 73. Split the log change in the futures price into spot and roll, and check
  \cref{eq4:roll} using the basis at entry and exit.
  \item\label{ex4:ten} Show that a European call and put with the same strike $K$ and expiry $T$ satisfy
  $C-P=(F_0-K)e^{-rT}$, where $F_0$ is the forward price, for any asset with a forward market and a
  constant rate $r$. Why does this make the put-call parity of the proof of \cref{thm4:crr} a special
  case?
\end{enumerate}

\subsection*{Answers}

\begin{enumerate}[label=\textbf{4.\arabic*}, leftmargin=*]
  \item $D(1)=1/1.03=\exOneDOne$. The bond prices at par, $0.035D(1)+1.035D(2)=1$, so $D(2)=\exOneDTwo$
  and $F(1,2)=D(1)/D(2)-1=\exOneFwd\%$.
  \item $F=50e^{(0.05-0.02)\times0.75}=\exTwoF$.
  \item By \cref{eq4:fwdvalue}, $V=55e^{-0.02/3}-52e^{-0.05/3}=\exThreeV$.
  \item $F=1.35\times1.015/1.025=\exFourF$. The quote is above parity, so sell US dollars forward:
  borrow 1,350,000 Canadian dollars, buy one million US dollars and invest them at 5\% to receive
  1,025,000. Selling these forward at 1.34 brings \exFourGet{} Canadian dollars against a loan balance
  of \exFourOwe, a profit of \exFourProfit.
  \item $\psi_u+\psi_d=1$ and $110\psi_u+90\psi_d=100$ give $\psi_u=\exFivePsiU$ and
  $\psi_d=\exFivePsiD$. The call is worth $10\psi_u=\exFiveCall$ and the put $15\psi_d=\exFivePut$.
  \item $\eta=r+\lambda-\ln(F/S)=0.06+\ln(70/66)=\exSixY\%$.
  \item $q=(1.02-0.9)/0.2=\exSevenQ$. The put pays 0 or 5, so its price is $0.4\times5/1.02=\exSevenP$
  and $\Delta=(0-5)/(55-45)=\exSevenDelta$: sell half a share.
  \item The nodes are 55 and 45 after one step and 60.5, 49.5 and 40.5 after two, where the put pays
  0, 2.5 and 11.5. With $q=0.6$, the continuation value at 55 is $(0.6\times0+0.4\times2.5)/1.02=\exEightContU$
  and exercise pays nothing, so the holder waits. At 45 the continuation value is
  $(0.6\times2.5+0.4\times11.5)/1.02=\exEightContD$ and exercise pays \exEightExD, so the put is
  exercised. The American put is worth $(0.6\times\exEightContU+0.4\times7)/1.02=\exEightAm$; at the root
  exercise pays 2.00, less than \exEightAm, so the holder waits. The European put is \exEightEu.
  \item In percent, the log futures change, the spot return and the roll are
  \begin{equation*}
    100\ln\tfrac{74}{72}=\exNineEx,\qquad 100\ln\tfrac{73}{70}=\exNineSpot,\qquad
    \exNineEx-\exNineSpot=\exNineRoll .
  \end{equation*}
  The basis was $100\ln(72/70)=\exNineBasisIn$ at entry and $100\ln(74/73)=\exNineBasisOut$ at exit,
  and the change in the basis from entry to exit is the same roll. The curve was in contango and the
  contract lost ground on the spot.
  \item The payoff of the call less the put is $(S_T-K)^+-(K-S_T)^+=S_T-K$, which is the payoff of a long
  forward with delivery price $K$. By \cref{eq4:fwdvalue} at $t=0$ that is worth $(F_0-K)e^{-rT}$, and
  the law of one price gives the result. For a stock without dividends $F_0=S_0e^{rT}$, which turns it
  into $C-P=S_0-Ke^{-rT}$.
\end{enumerate}

% =====================================================================================
\section{Notes, sources and credits}
\label{sec4:credits}

\paragraph{Sources}
\textcite{ross1978} gave the one-period fundamental theorem. \textcite{harrison1979} tied the
absence of arbitrage to a martingale measure through the viability of prices, and
\textcite{harrison1981} proved the exact equivalence for finitely many states and dates and developed
the theory of continuous trading. The proof by separation given
here follows the textbook treatments of \textcite{follmer2016} and \textcite{duffie2001}, and the
open interval of prices for a claim that cannot be replicated follows \textcite{follmer2016}. The relation between futures and forward prices is from \textcite{cox1981}, the short rate
model from \textcite{vasicek1977}, and the commodity model from \textcite{schwartz1997}. The convenience yield is from \textcite{kaldor1939} and normal backwardation
from \textcite{keynes1930}; \textcite{gorton2006} and \textcite{erb2006} study the returns of rolled
commodity futures, and \textcite{du2018} the failure of covered interest parity. The binomial model is
from \textcite{cox1979}, the refined tree from \textcite{leisen1996}, using the normal approximation
of \textcite{peizer1968}, and the error expansion from \textcite{diener2004}. The results on early
exercise go back to \textcite{merton1973}. The case study follows \textcite{fama1987} with the
standard errors of \textcite{newey1987}. As standard references, \textcite{shreve2004a} develops the
binomial model and \textcite{shreve2004} the continuous-time theory, including the perpetual put, and
\textcite{hull2022} covers market conventions for forwards, futures and currency quotes.

\paragraph{Open source projects}
Software used to produce the results, with the versions installed in the recorded environment: the
chapter's scripts use NumPy 2.2.6 \parencite{numpy} for arrays and random draws, SciPy 1.18.1 \parencite{scipy},
including its HiGHS linear programming solver for the bounds and arbitrage tests, statsmodels 0.15.0
\parencite{statsmodels} for the regressions and Matplotlib 3.11.2 \parencite{matplotlib} for the
figures. The cross-checks use base R 4.5.0 \parencite{rcore} and GNU Octave 9.4.0 with its built-in GLPK
solver \parencite{octave}. Projects cited for comparison, none of which the chapter calls: QuantLib
\parencite{quantlib} and FinancePy \parencite{financepy} implement curve bootstrapping, forward and
futures pricing and binomial engines, including the Leisen-Reimer tree, and the notebooks of
\textcite{cantaro86} and \textcite{py4fi} treat binomial trees and their convergence. The full list,
with links, is in \texttt{CREDITS.md}.

\paragraph{Data}
Treasury bill and benchmark bond yields are from the Bank of Canada Valet service
\parencite{bocvalet}. The bill series are \texttt{TB.CDN.90D.MID}, \texttt{TB.CDN.180D.MID} and
\texttt{TB.CDN.1Y.MID}. The bond series are \texttt{BD.CDN.$x$.DQ.YLD} with $x$ equal to
\texttt{2YR}, \texttt{3YR}, \texttt{5YR}, \texttt{7YR}, \texttt{10YR} and \texttt{LONG}. The saved file
holds daily observations from January 2001 to 29 September 2026, subject to each series' own gaps
and non-publication dates, retrieved on 19 October 2026. The long bond series is saved with the others
but not used, since the curve stops at ten years. The
download script is \texttt{scripts/pull\_boc.py}; the saved file is \texttt{data/goc\_bills\_bonds.csv}.
Bill yields
are published mid rates and bond yields are benchmark yields to maturity, both in percent; the chapter
treats them as described in \cref{sec4:curve}. The Bank can revise published values, and the numbers
here use the saved file.

\printbibliography[heading=subbibliography,title={References for Chapter \thechapter}]
\end{refsection}
